\documentclass[12pt,a4paper,UTF8]{ctexart}
\usepackage[a4paper,top=22mm,bottom=22mm,left=22mm,right=22mm,headheight=15pt]{geometry}
\usepackage{amsmath,amssymb,mathtools,bm}
\usepackage{graphicx}
\usepackage{xcolor}
\usepackage{booktabs,longtable,array}
\usepackage{enumitem}
\usepackage{fancyhdr}
\usepackage{caption}
\usepackage{float}
\usepackage{microtype}
\usepackage{setspace}
\usepackage[hidelinks,unicode]{hyperref}
\usepackage[most]{tcolorbox}

\definecolor{ContestBlue}{HTML}{17365D}
\definecolor{SoftBlue}{HTML}{EEF4FA}
\definecolor{SoftGray}{HTML}{F5F6F7}
\hypersetup{pdftitle={第 091 套},pdfauthor={}}

\setstretch{1.24}
\setlength{\parindent}{2em}
\setlength{\parskip}{0.35em}
\setlist{itemsep=0.35em,topsep=0.4em}
\captionsetup{font=small,labelfont=bf,justification=centering}
\setcounter{secnumdepth}{0}

\pagestyle{fancy}
\fancyhf{}
\fancyhead[L]{\small 第 091 套}
\fancyfoot[C]{\thepage}
\renewcommand{\headrulewidth}{0.4pt}

\ctexset{
  section={format=\Large\bfseries\color{ContestBlue},beforeskip=1.2em,afterskip=0.7em},
  subsection={format=\large\bfseries\color{ContestBlue},beforeskip=1.0em,afterskip=0.5em},
  subsubsection={format=\normalsize\bfseries,beforeskip=0.8em,afterskip=0.35em}
}

\makeatletter
\newcommand{\tightlist}{\setlength{\itemsep}{0pt}\setlength{\parskip}{0pt}}
\newcommand{\maxwidth}{\ifdim\Gin@nat@width>\linewidth\linewidth\else\Gin@nat@width\fi}
\newcommand{\maxheight}{\ifdim\Gin@nat@height>0.72\textheight 0.72\textheight\else\Gin@nat@height\fi}
\newcommand{\pandocbounded}[1]{\begingroup\setkeys{Gin}{width=\maxwidth,height=\maxheight,keepaspectratio}#1\endgroup}
\makeatother

\begin{document}

\begin{titlepage}
  \centering
  \vspace*{58mm}
  {\Huge\bfseries\color{ContestBlue} 第 091 套\par}
  \vfill
  \vspace*{28mm}
\end{titlepage}


\section{第 1 题}\label{ux7b2c-1-ux9898}

两面墙之间水平固定一根光滑杆，杆上套有一个质量为 \(m\)
的小环。小环通过一根长为 \(l\)、不可伸长且质量不计的细线连接一个质量为
\(2m\)
的小球。小球可在杆与竖直方向所确定的平面内摆动，两面墙之间的距离足够大。

\begin{enumerate}
\def\labelenumi{\arabic{enumi}.}
\tightlist
\item
  若小球偏离竖直方向的最大角度为
  \(30^\circ\)，求小球到达最大摆角时小环的加速度大小以及小球相对地面的加速度大小。
\item
  若摆角很小，求系统的振动周期。
\end{enumerate}

\begin{figure}[H]
\centering
\includegraphics[width=0.9\linewidth,height=\textheight,keepaspectratio,alt={图 1-1　可沿水平光滑杆移动的悬点与摆球系统。}]{../figures/question-01.pdf}
\caption*{图 1-1　可沿水平光滑杆移动的悬点与摆球系统。}
\end{figure}

\subsection{一、题型定位}\label{ux4e00ux9898ux578bux5b9aux4f4d}

小环是可动悬点，摆球的水平张力会推动它。对小环和摆球组成的系统写动量关系，再配合摆球径向方程，才能求有限摆角处的加速度。

\subsection{二、知识点}\label{ux4e8cux77e5ux8bc6ux70b9}

\begin{itemize}
\tightlist
\item
  可动悬点单摆的相对运动与牵连运动。
\item
  最大摆角处 \(v=0\) 与 \(a\ne0\) 的区别。
\item
  随小环一起平动的非惯性系及惯性力。
\item
  细线方向、切线方向的分解。
\item
  水平方向无外力时的质心运动守恒。
\item
  小角近似 \(\sin\theta\simeq\theta\)、\(\cos\theta\simeq1\)。
\item
  二自由度系统约束消元与等效角频率。
\item
  加速度的矢量合成及夹角判断。
\end{itemize}

\subsection{三、物理图像}\label{ux4e09ux7269ux7406ux56feux50cf}

设小球位于小环左下方，细线与竖直线的夹角为
\(\theta\)。细线对小环的拉力具有向左的水平分量，所以小环向左加速。若改在随小环平动的参考系中观察摆球，就必须给摆球附加一个水平向右的惯性力。这个惯性力既有沿细线的分量，又有沿摆球轨迹切线的分量。

在最大摆角处，摆球相对于小环的速度为零，因此相对法向加速度
\(v_{\rm rel}^2/l\)
为零；但重力和惯性力在切线方向的合力不为零，摆球会立刻向平衡位置加速。

小振动时，摆球向左偏，整个系统的水平质心仍不能加速。为了抵消摆球质心的左移，小环也必须随摆角变化而移动。于是摆球相对于悬点的有效惯性不同于固定悬点单摆，振动周期也不再是
\(2\pi\sqrt{l/g}\)。

\subsection{四、解题入口}\label{ux56dbux89e3ux9898ux5165ux53e3}

第一问选择随小环一起平动、但不随细线转动的参考系。设小环加速度大小为
\(a\)，摆球在该参考系中受到大小为

\[F^*=2ma\]

的惯性力。先由小环的水平动力学把 \(a\) 与张力 \(T\)
联系起来，再在摆球的径向方程中求出 \(T\)，最后写切向方程求相对加速度。

第二问可从水平质心关系直接切入。令小环向右的位置为
\(x\)，并规定摆球向左偏时 \(\theta>0\)，则摆球横坐标为

\[x_b=x-l\sin\theta.\]

系统水平方向不受外力，所以

\[m\ddot x+2m\ddot x_b=0.\]

将此式线性化，再与摆球相对于悬点的切向方程联立，就能得到单一的
\(\theta\) 方程。

\subsection{五、解析}\label{ux4e94ux89e3ux6790}

图 1-2
左侧给出最大摆角处摆球的等效受力与相对切向加速度，右侧给出小振动时小环加速度
\(a\)、摆球切向加速度 \(l\beta\)
及二者的约束关系。下文第（1）---（3）步使用左图，第（4）---（5）步使用右图。

\begin{figure}[H]
\centering
\includegraphics[width=0.9\linewidth,height=\textheight,keepaspectratio,alt={图 1-2　最大摆角处的受力分解与小振动约束关系。}]{../figures/solution-01.pdf}
\caption*{图 1-2　最大摆角处的受力分解与小振动约束关系。}
\end{figure}

\textbf{（1）最大摆角处小环的加速度} 先看图 1-2
左侧。最大摆角处摆球相对小环的径向速度和径向加速度均为零，图中的
\(T\)、\(2mg\) 与随小环参考系引入的惯性力 \(F^*=2ma\) 正好用于径向投影。

记最大摆角为

\[\theta_0=30^\circ.\]

先对质量为 \(m\)
的小环列水平方向牛顿第二定律。小环在水平方向只受到细线张力的水平分量，因此

\[ma=T\sin\theta_0.\]

于是

\[
\boxed{a=\frac{T\sin\theta_0}{m}}.
\]

为了处理摆球，转到随小环平动的参考系。该参考系以加速度 \(a\)
向左运动，所以质量为 \(2m\) 的摆球受到向右的惯性力

\[
F^*=2ma.
\]

利用小环方程 \(ma=T\sin\theta_0\)，可以把惯性力改写成

\[
F^*=2T\sin\theta_0.
\]

在最大摆角处，摆球相对于小环的瞬时速度为零，所以相对径向加速度为

\[
a_{r,\rm rel}=\frac{v_{\rm rel}^2}{l}=0.
\]

现在沿``从摆球指向小环''的径向正方向投影。张力 \(T\)
沿径向正方向；向右的惯性力与径向的夹角为
\(90^\circ-\theta_0\)，故其径向分量是
\(F^*\sin\theta_0\)；重力的径向分量与正方向相反，大小为
\(2mg\cos\theta_0\)。因此

\[
T+F^*\sin\theta_0-2mg\cos\theta_0=0.
\]

代入 \(F^*=2T\sin\theta_0\)：

\[
T+2T\sin^2\theta_0=2mg\cos\theta_0.
\]

再代入 \(\sin30^\circ=1/2\)、\(\cos30^\circ=\sqrt3/2\)：

\[
\begin{aligned}
T+2T\left(\frac12\right)^2
&=2mg\frac{\sqrt3}{2},\\
T+\frac12T
&=\sqrt3mg,\\
\frac32T
&=\sqrt3mg.
\end{aligned}
\]

所以

\[
\boxed{T=\frac{2\sqrt3}{3}mg}.
\]

回代小环方程：

\[
\begin{aligned}
a
&=\frac{T\sin30^\circ}{m}\\
&=\frac{1}{m}\cdot\frac{2\sqrt3}{3}mg\cdot\frac12\\
&=\frac{\sqrt3}{3}g.
\end{aligned}
\]

故小环的加速度大小为

\[
\boxed{a_{\text{环}}=\frac{g}{\sqrt3}}.
\]

这里必须特别核对径向投影：惯性力的径向分量是 \(F^*\sin30^\circ\)，不是
\(F^*\cos30^\circ\)。如果投影函数写错，恰好会把后续张力和加速度全部带偏。

\textbf{（2）最大摆角处摆球相对于小环的切向加速度}

仍沿用左图绿色箭头的方向，取摆球回到最低点的切线方向为正方向。重力的切向分量为

\[2mg\sin\theta_0,\]

向右惯性力的切向分量为

\[F^*\cos\theta_0.\]

二者在图示位置都指向平衡位置，故相加：

\[
2m a_{\rm rel}=2mg\sin\theta_0+F^*\cos\theta_0.
\]

在 \(\theta_0=30^\circ\) 时，已经有

\[F^*=2T\sin30^\circ=T=\frac{2\sqrt3}{3}mg.\]

于是

\[
\begin{aligned}
2m a_{\rm rel}
&=2mg\cdot\frac12
+\frac{2\sqrt3}{3}mg\cdot\frac{\sqrt3}{2}\\
&=mg+mg\\
&=2mg.
\end{aligned}
\]

因此

\[
\boxed{a_{\rm rel}=g}.
\]

\textbf{（3）把相对加速度还原为摆球对地加速度}

按图 1-2 左侧红、绿两箭头，小环加速度水平向左，大小为
\(g/\sqrt3\)；摆球相对小环的切向加速度沿切线向右下方，大小为
\(g\)。细线与竖直方向夹角为 \(30^\circ\)，所以切线与水平向右方向夹角也是
\(30^\circ\)。

取水平向右为 \(x\) 正方向、竖直向上为 \(y\) 正方向。小环加速度为

\[
\bm a_{\text{环}}=\left(-\frac{g}{\sqrt3},0\right).
\]

摆球相对小环的切向加速度为

\[
\bm a_{\rm rel}
=g\bigl(\cos30^\circ,-\sin30^\circ\bigr)
=\left(\frac{\sqrt3}{2}g,-\frac12g\right).
\]

因为小环参考系只做平动，没有转动，摆球对地加速度就是牵连加速度与相对加速度之和：

\[
\begin{aligned}
\bm a_b
&=\bm a_{\text{环}}+\bm a_{\rm rel}\\
&=\left(-\frac{g}{\sqrt3}+\frac{\sqrt3}{2}g,-\frac12g\right)\\
&=\left(\frac{\sqrt3}{6}g,-\frac12g\right).
\end{aligned}
\]

故其大小为

\[
\begin{aligned}
a_b
&=\sqrt{\left(\frac{\sqrt3}{6}g\right)^2+
\left(-\frac12g\right)^2}\\
&=g\sqrt{\frac{3}{36}+\frac14}\\
&=g\sqrt{\frac1{12}+\frac3{12}}\\
&=\frac{g}{\sqrt3}.
\end{aligned}
\]

所以摆球相对地面的加速度大小为

\[
\boxed{a_{\text{球}}=\frac{g}{\sqrt3}}.
\]

两者大小相同只是本题质量比 \(m:2m\) 与最大摆角 \(30^\circ\)
共同造成的结果，方向并不相同。

\textbf{（4）小振动方程：建立水平质心约束}

转到图 1-2 右侧的小振动示意。令小环横坐标为 \(x\)，摆球向左偏时
\(\theta>0\)。图中水平箭头 \(a\) 表示小环加速度，斜向箭头 \(l\beta\)
表示摆球相对小环的切向加速度。摆球横坐标为

\[x_b=x-l\sin\theta.\]

整个系统水平方向没有外力，因此

\[
m\ddot x+2m\ddot x_b=0.
\]

先对 \(x_b\) 求两次导数：

\[
\dot x_b=\dot x-l\cos\theta\,\dot\theta,
\]

\[
\ddot x_b
=\ddot x+l\sin\theta\,\dot\theta^2-l\cos\theta\,\ddot\theta.
\]

小振动中 \(\theta\) 很小，\(\dot\theta^2\sin\theta\) 是高阶小量，且
\(\cos\theta\simeq1\)，故

\[
\ddot x_b\simeq\ddot x-l\ddot\theta.
\]

代入质心约束：

\[
m\ddot x+2m(\ddot x-l\ddot\theta)=0.
\]

展开并约去 \(m\)：

\[
3\ddot x-2l\ddot\theta=0.
\]

因此

\[
\boxed{\ddot x=\frac23l\ddot\theta}.
\]

\textbf{（5）小振动方程：写出摆角动力学}

可直接由拉格朗日方程得到摆角方程。系统动能为

\[
\begin{aligned}
K
&=\frac12m\dot x^2
+\frac12(2m)\left[
(\dot x-l\cos\theta\dot\theta)^2
(l\sin\theta\dot\theta)^2\right]\\
&=\frac32m\dot x^2
-2ml\dot x\cos\theta\dot\theta
+ml^2\dot\theta^2.
\end{aligned}
\]

以最低点为零势能面，摆球势能为

\[U=2mgl(1-\cos\theta).\]

在小角近似下，关于 \(\theta\) 的线性化方程为

\[
l\ddot\theta-\ddot x+g\theta=0.
\]

它也可由随小环参考系的切向牛顿方程得到：相对切向加速度为
\(l\ddot\theta\)，重力提供 \(-g\theta\)
的回复项，小环水平加速度在切向上贡献 \(\ddot x\)。

代入 \(\ddot x=(2/3)l\ddot\theta\)：

\[
l\ddot\theta-\frac23l\ddot\theta+g\theta=0.
\]

合并同类项：

\[
\frac13l\ddot\theta+g\theta=0.
\]

两边乘以 \(3/l\)：

\[
\boxed{\ddot\theta+\frac{3g}{l}\theta=0}.
\]

因此系统角频率为

\[
\omega_0=\sqrt{\frac{3g}{l}},
\]

振动周期为

\[
\boxed{T=2\pi\sqrt{\frac{l}{3g}}}.
\]

\textbf{另一种核验：约化质量法}

由线性化的水平质心约束可积分得到

\[3\dot x-2l\dot\theta=0,\]

若系统初始总水平动量为零，则

\[\dot x=\frac23l\dot\theta.\]

将其代入小角动能

\[K\simeq\frac32m\dot x^2-2ml\dot x\dot\theta+ml^2\dot\theta^2,\]

逐项化简：

\[
\begin{aligned}
K
&=\frac32m\left(\frac23l\dot\theta\right)^2
-2ml\left(\frac23l\dot\theta\right)\dot\theta
+ml^2\dot\theta^2\\
&=\frac23ml^2\dot\theta^2-\frac43ml^2\dot\theta^2+ml^2\dot\theta^2\\
&=\frac13ml^2\dot\theta^2.
\end{aligned}
\]

把它写成 \(K=\frac12I_{\rm eff}\dot\theta^2\)，得到

\[I_{\rm eff}=\frac23ml^2.\]

小角势能

\[U\simeq mgl\theta^2=\frac12(2mgl)\theta^2.\]

故

\[\omega_0^2=\frac{2mgl}{(2/3)ml^2}=\frac{3g}{l},\]

与前述结果一致。这一核验能帮助学生理解：悬点可动后，系统对摆角的等效转动惯量减小，所以角频率增大、周期缩短。

\newpage

\section{第 2 题}\label{ux7b2c-2-ux9898}

一条均匀、柔软且不可伸长的长链可视作``贪吃蛇''，总长为 \(L\)、总质量为
\(m\)。长链在前端附近平滑地折回：后端沿直线以速度 \(v\)
运动，前端沿相反方向以恒定速度 \(u\)
被拉动。忽略重力以及弯折区的尺寸。为了维持这一稳定运动，作用在前端的拉力
\(F\) 应为多大？

\begin{figure}[H]
\centering
\includegraphics[width=0.9\linewidth,height=\textheight,keepaspectratio,alt={图 2-1　均匀柔软长链在移动弯折处改变运动方向。}]{../figures/question-02.pdf}
\caption*{图 2-1　均匀柔软长链在移动弯折处改变运动方向。}
\end{figure}

\subsection{一、题型定位}\label{ux4e00ux9898ux578bux5b9aux4f4d-1}

弯折处不断把链条动量从一个方向转到相反方向。对随弯折点移动的控制体写动量流，外力大小由质量流率与速度变化共同决定。

\subsection{二、知识点}\label{ux4e8cux77e5ux8bc6ux70b9-1}

\begin{itemize}
\tightlist
\item
  线密度 \(\lambda=m/L\)。
\item
  不可伸长柔索的长度约束。
\item
  移动控制面与质量通量。
\item
  一维动量的代数符号。
\item
  变质量系统应使用动量流，而不能直接套整体 \(F=ma_{\rm cm}\)
  而不求质心加速度。
\item
  极限检验：\(u=v\)、\(v=0\) 等情形。
\item
  功率与动能通量的交叉核验。
\end{itemize}

\subsection{三、物理图像}\label{ux4e09ux7269ux7406ux56feux50cf-1}

把向右规定为正方向。上支链条以速度 \(+u\) 运动，下支链条以速度 \(-v\)
运动。弯折区向右移动时，上支长度的增加量与下支长度的减少量并不分别等于
\(u\,\mathrm{d}t\) 或 \(v\,\mathrm{d}t\)，因为弯折点本身也在移动。

在时间 \(\mathrm{d}t\) 内，某一小段链从下支的速度 \(-v\)
过渡到上支的速度 \(+u\)。每单位质量的水平动量改变量为
\(u-(-v)=u+v\)。与此同时，穿过弯折区的质量流率也与 \(u+v\)
成正比，所以最终拉力含有 \((u+v)^2\)。

\subsection{四、解题入口}\label{ux56dbux89e3ux9898ux5165ux53e3-1}

设弯折点在 \(\mathrm{d}t\) 内向右移动 \(\mathrm{d}X\)。令上支长度变化为
\(\mathrm{d}L_1\)，下支长度变化为
\(\mathrm{d}L_2\)。从端点运动与弯折点运动的相对关系写出

\[
\mathrm{d}L_1=u\,\mathrm{d}t-\mathrm{d}X,
\]

\[
\mathrm{d}L_2=-v\,\mathrm{d}t-\mathrm{d}X.
\]

总长度不变，故 \(\mathrm{d}L_1+\mathrm{d}L_2=0\)。由此先求
\(\mathrm{d}X\)、\(\mathrm{d}L_1\) 与
\(\mathrm{d}L_2\)，再把两条直线段的水平动量写成

\[p=\lambda L_1u-\lambda L_2v.\]

\subsection{五、解析}\label{ux4e94ux89e3ux6790-1}

图 2-2 把弯折区包在移动控制体内：上支速度为 \(+u\)，下支速度为
\(-v\)，右端拉力为 \(F\)。图下方分别列出弯折点位移
\(\mathrm dX\)、真正转入上支的链长 \(\mathrm d\ell\) 与动量增量
\(\mathrm dp\)；三者不可混为同一个长度变化。

\begin{figure}[H]
\centering
\includegraphics[width=0.865\linewidth,height=\textheight,keepaspectratio,alt={图 2-2　随弯折区移动的控制体及长度、动量通量。}]{../figures/solution-02.pdf}
\caption*{图 2-2　随弯折区移动的控制体及长度、动量通量。}
\end{figure}

\textbf{（1）确定线密度与正方向} 沿图 2-2
的水平向右方向取正。长链均匀，线密度为

\[
\boxed{\lambda=\frac{m}{L}}.
\]

规定前端被拉动的方向为水平向右正方向。因此上支速度为 \(+u\)，下支速度为
\(-v\)。

设初始时上、下两直线段的长度分别为 \(L_1\)、\(L_2\)。弯折区长度忽略时

\[L_1+L_2=L.\]

\textbf{（2）由不可伸长条件求弯折点位移}

在短时间 \(\mathrm{d}t\) 内，设图中弯折点向右移动
\(\mathrm{d}X\)。这里先求控制体边界的移动量，再求穿过弯折区的链长通量。

对上支而言，右端以 \(u\) 向右运动会使长度增加
\(u\,\mathrm{d}t\)，但弯折点也向右移动
\(\mathrm{d}X\)，会抵消一部分长度增加。因此

\[
\mathrm{d}L_1=u\,\mathrm{d}t-\mathrm{d}X.
\]

对下支而言，左端以速度 \(v\) 向左运动，使其右向坐标减少
\(v\,\mathrm{d}t\)；弯折点又向右移动
\(\mathrm{d}X\)，二者都使下支长度减小，所以

\[
\mathrm{d}L_2=-v\,\mathrm{d}t-\mathrm{d}X.
\]

链条不可伸长，总长度保持不变：

\[\mathrm{d}L_1+\mathrm{d}L_2=0.\]

代入：

\[
u\,\mathrm{d}t-\mathrm{d}X-v\,\mathrm{d}t-\mathrm{d}X=0.
\]

整理得

\[
2\,\mathrm{d}X=(u-v)\,\mathrm{d}t,
\]

所以弯折点位移为

\[
\boxed{\mathrm{d}X=\frac12(u-v)\,\mathrm{d}t}.
\]

于是弯折点速度为

\[
\boxed{V_f=\frac{\mathrm{d}X}{\mathrm{d}t}=\frac{u-v}{2}}.
\]

将 \(\mathrm{d}X\) 回代到 \(\mathrm{d}L_1\)：

\[
\begin{aligned}
\mathrm{d}L_1
&=u\,\mathrm{d}t-\frac12(u-v)\,\mathrm{d}t\\
&=\frac12(u+v)\,\mathrm{d}t.
\end{aligned}
\]

同理

\[
\begin{aligned}
\mathrm{d}L_2
&=-v\,\mathrm{d}t-\frac12(u-v)\,\mathrm{d}t\\
&=-\frac12(u+v)\,\mathrm{d}t.
\end{aligned}
\]

因此

\[
\boxed{\mathrm{d}L_1=-\mathrm{d}L_2=\frac12(u+v)\,\mathrm{d}t}.
\]

这说明在 \(\mathrm{d}t\) 内有长度

\[\mathrm{d}\ell=\frac12(u+v)\,\mathrm{d}t\]

的链元从下支转入上支。

\textbf{（3）计算全链动量的变化}

图 2-2 中同一段链元从下支的速度 \(-v\) 变成上支的速度
\(+u\)。弯折区质量忽略时，初始水平动量为

\[
p_i=\lambda L_1u-\lambda L_2v.
\]

经过 \(\mathrm{d}t\) 后，上支长度变成
\(L_1+\mathrm{d}L_1\)，下支长度变成 \(L_2+\mathrm{d}L_2\)，故

\[
p_f=\lambda(L_1+\mathrm{d}L_1)u-\lambda(L_2+\mathrm{d}L_2)v.
\]

两式相减：

\[
\begin{aligned}
\mathrm{d}p
&=p_f-p_i\\
&=\lambda u\,\mathrm{d}L_1-\lambda v\,\mathrm{d}L_2.
\end{aligned}
\]

由于 \(\mathrm{d}L_2=-\mathrm{d}L_1\)，所以

\[
\mathrm{d}p=\lambda(u+v)\,\mathrm{d}L_1.
\]

再代入 \(\mathrm{d}L_1=\frac12(u+v)\,\mathrm{d}t\)：

\[
\begin{aligned}
\mathrm{d}p
&=\lambda(u+v)\cdot\frac12(u+v)\,\mathrm{d}t\\
&=\frac{\lambda}{2}(u+v)^2\,\mathrm{d}t.
\end{aligned}
\]

因此全链水平动量的变化率为

\[
\frac{\mathrm{d}p}{\mathrm{d}t}=\frac{\lambda}{2}(u+v)^2.
\]

按题目模型，维持运动的水平外力即前端拉力 \(F\)，故

\[
F=\frac{\mathrm{d}p}{\mathrm{d}t}.
\]

所以

\[
\boxed{F=\frac{\lambda}{2}(u+v)^2}.
\]

最后代入 \(\lambda=m/L\)：

\[
\boxed{F=\frac{m}{2L}(u+v)^2}.
\]

\textbf{（4）从单个链元看同一结果}

单位时间穿过弯折区的质量为

\[
\dot m_f=\lambda\frac{\mathrm{d}\ell}{\mathrm{d}t}
=\frac{\lambda}{2}(u+v).
\]

每单位质量的水平速度由 \(-v\) 变为 \(+u\)，所以速度改变量为

\[\Delta v_x=u-(-v)=u+v.\]

动量流率即

\[
\dot p_x=\dot m_f\Delta v_x
=\frac{\lambda}{2}(u+v)^2,
\]

再次得到

\[F=\frac{m}{2L}(u+v)^2.\]

\textbf{功率核验及其边界}

若只计算链条平直部分的动能

\[K=\frac12\lambda L_1u^2+\frac12\lambda L_2v^2,\]

则

\[
\begin{aligned}
\frac{\mathrm{d}K}{\mathrm{d}t}
&=\frac12\lambda u^2\frac{\mathrm{d}L_1}{\mathrm{d}t}
+\frac12\lambda v^2\frac{\mathrm{d}L_2}{\mathrm{d}t}\\
&=\frac14\lambda(u+v)(u^2-v^2)\\
&=\frac14\lambda(u-v)(u+v)^2.
\end{aligned}
\]

这个功率一般不等于
\(Fu\)，因为题目同时规定了后端的运动，后端约束也可能交换功；弯折区还可能承受维持其形状的约束力。因而本题最可靠的是整体动量变化，而不是只凭前端功率单独反推
\(F\)。课堂上应借此强调：在多端口连续介质问题中，功率核验必须把所有外力作用点都列全。

\textbf{极限核验}

当 \(u=v=U\) 时，弯折点静止，单位时间有质量流率 \(\lambda U\)
的链元由速度 \(-U\) 变为 \(+U\)，故

\[F=\lambda U(2U)=2\lambda U^2.\]

一般式 \(F=\lambda(u+v)^2/2\) 也给出 \(2\lambda U^2\)，与这个特例一致。

\newpage

\section{第 3 题}\label{ux7b2c-3-ux9898}

两根电阻均为 \(r\)、质量分别为 \(m\) 和 \(2m\)
的光滑导体棒，垂直放置在间距为 \(l\)
的两条水平光滑导轨之间。整个装置处在垂直纸面向里、磁感应强度为 \(B\)
的匀强磁场中。两棒所在支路各接有一个电容均为 \(C\)
的电容器。初始时两棒静止，两电容器均不带电。自 \(t=0\) 起，对质量为
\(2m\) 的右棒施加水平外力

\[F(t)=F_0\cos\omega t.\]

\begin{enumerate}
\def\labelenumi{\arabic{enumi}.}
\tightlist
\item
  求两棒质心速度 \(v_C(t)\)。
\item
  求回路电动势 \(\mathcal E(t)\)。
\item
  求右棒、左棒的速度 \(v_1(t)\)、\(v_2(t)\)。
\end{enumerate}

\begin{figure}[H]
\centering
\includegraphics[width=0.9\linewidth,height=\textheight,keepaspectratio,alt={图 3-1　匀强磁场中的双导体棒、电容支路及简谐外力。}]{../figures/question-03.pdf}
\caption*{图 3-1　匀强磁场中的双导体棒、电容支路及简谐外力。}
\end{figure}

\subsection{一、题型定位}\label{ux4e00ux9898ux578bux5b9aux4f4d-2}

外力驱动两根导体棒，棒速产生电动势，电流又反过来改变棒的运动。先消去两棒共有的动量变量，再解受迫电路方程，可把机械与电学部分分开。

\subsection{二、知识点}\label{ux4e8cux77e5ux8bc6ux70b9-2}

\begin{itemize}
\tightlist
\item
  系统总动量定理与质心速度。
\item
  动生电动势 \(\mathcal E=Blv\)。
\item
  安培力 \(F_B=Bil\) 及两棒受力方向。
\item
  电流与电荷的关系 \(i=\dot q\)。
\item
  电容电压 \(U_C=q/C\)。
\item
  基尔霍夫电压定律。
\item
  一阶常系数非齐次微分方程。
\item
  简谐特解与指数暂态。
\item
  初始条件 \(q(0)=0\)、\(\dot q(0)=0\) 的一致性检查。
\end{itemize}

\subsection{三、物理图像}\label{ux4e09ux7269ux7406ux56feux50cf-2}

右棒受外力向右运动时产生动生电动势。若两棒速度不同，两棒的动生电动势不能完全抵消，回路中便有电流。电流流过两根棒时，安培力方向相反：它们是系统内部的相互作用，能改变两棒的相对运动，却不能改变两棒总动量。

电容使电荷 \(q\) 成为系统的状态变量，电阻使暂态按 \(e^{-Dt/(2r)}\)
衰减。因为本题没有电感，回路方程对电流只含一阶导数；把速度消去后得到的仍是关于
\(q\) 的一阶方程，而不是电磁振荡方程。

\subsection{四、解题入口}\label{ux56dbux89e3ux9898ux5165ux53e3-2}

规定两棒向右速度为正，右棒质量为 \(2m\)、速度为 \(v_1\)，左棒质量为
\(m\)、速度为 \(v_2\)。取电流正方向使右棒安培力向右，则

\[
2m\dot v_1=Bl\dot q+F_0\cos\omega t,
\]

\[
m\dot v_2=-Bl\dot q.
\]

两式相加时安培力抵消，立即得到质心运动。回路方面选定

\[
\mathcal E(t)\equiv\mathcal E_2-\mathcal E_1,
\]

则基尔霍夫方程可写成

\[
\mathcal E(t)=2r\dot q+\frac{2q}{C}.
\]

再利用 \(\mathcal E_1=Blv_1\)、\(\mathcal E_2=Blv_2\)，便可消去
\(v_1-v_2\)。

\subsection{五、解析}\label{ux4e94ux89e3ux6790-2}

图 3-2 左侧是按电流正方向 \(i=\dot q\)
画出的等效回路，右侧依次给出两棒动力学、回路方程和消元后的电荷方程。下文第（2）步对应右上框，第（3）步对应左图和右中框，第（4）步完成右上框到右下框的消元。

\begin{figure}[H]
\centering
\includegraphics[width=0.9\linewidth,height=\textheight,keepaspectratio,alt={图 3-2　等效回路、两棒动力学与电荷方程的消元关系。}]{../figures/solution-03.pdf}
\caption*{图 3-2　等效回路、两棒动力学与电荷方程的消元关系。}
\end{figure}

\textbf{（1）质心速度} 把两棒看作一个整体。水平方向外力只有
\(F_0\cos\omega t\)，安培力是两棒之间经由电路传递的内部作用，总和为零。因此

\[
(2m+m)\dot v_C=F_0\cos\omega t.
\]

即

\[
3m\dot v_C=F_0\cos\omega t.
\]

初始时两棒静止，所以 \(v_C(0)=0\)。从 \(0\) 到 \(t\) 积分：

\[
\begin{aligned}
v_C(t)
&=\int_0^t\frac{F_0}{3m}\cos\omega t'\,\mathrm{d}t'\\
&=\frac{F_0}{3m\omega}\sin\omega t.
\end{aligned}
\]

故

\[
\boxed{v_C(t)=\frac{F_0}{3m\omega}\sin\omega t}.
\]

同时有总动量关系

\[
\boxed{2mv_1+mv_2=\frac{F_0}{\omega}\sin\omega t}.
\]

\textbf{（2）建立两棒动力学方程}

令 \(i=\dot q\)。按照图 3-2 左侧箭头选定的回路正方向，右棒受到向右安培力
\(Bil=Bl\dot q\)，左棒受到向左安培力；这正对应右上框中的两行方程。因此

\[
\boxed{2m\dot v_1=Bl\dot q+F_0\cos\omega t},
\]

\[
\boxed{m\dot v_2=-Bl\dot q}.
\]

由第二式并利用 \(v_2(0)=0\)、\(q(0)=0\)，直接积分：

\[
mv_2=-Blq,
\]

所以

\[
\boxed{v_2(t)=-\frac{Bl}{m}q(t)}.
\]

\textbf{（3）建立回路方程}

沿图 3-2
左侧回路从左棒开始绕行，依次经过左电容、两棒总电阻和右电容。两棒的动生电动势大小分别为

\[
\mathcal E_1=Blv_1,
\qquad
\mathcal E_2=Blv_2.
\]

按照选定的回路正方向，基尔霍夫电压定律为

\[
\mathcal E_1+\frac qC+r\dot q+r\dot q+\frac qC-\mathcal E_2=0.
\]

合并同类项：

\[
\frac{2q}{C}+2r\dot q+Bl(v_1-v_2)=0.
\]

也就是

\[
\boxed{2r\dot q+\frac{2q}{C}=Bl(v_2-v_1)}.
\]

因此，若定义回路净电动势为 \(\mathcal E=\mathcal E_2-\mathcal E_1\)，则

\[
\boxed{\mathcal E(t)=2r\dot q+\frac{2q}{C}}.
\]

若改用相反绕行方向，\(\mathcal E\)、\(q\)、\(i\)
的符号会同时反转，物理结果不变。

\textbf{（4）消去棒速，得到电荷方程}

现在把图 3-2 右上框中的 \(\dot v_1-\dot v_2\) 代入右中框。先对回路方程

\[
2r\dot q+\frac{2q}{C}+Bl(v_1-v_2)=0
\]

对时间求导：

\[
2r\ddot q+\frac{2}{C}\dot q+Bl(\dot v_1-\dot v_2)=0.
\]

由两棒动力学方程，

\[
\dot v_1=\frac{Bl}{2m}\dot q+\frac{F_0}{2m}\cos\omega t,
\]

\[
\dot v_2=-\frac{Bl}{m}\dot q.
\]

所以

\[
\begin{aligned}
\dot v_1-\dot v_2
&=\frac{Bl}{2m}\dot q+\frac{F_0}{2m}\cos\omega t
+\frac{Bl}{m}\dot q\\
&=\frac{3Bl}{2m}\dot q+\frac{F_0}{2m}\cos\omega t.
\end{aligned}
\]

代回：

\[
2r\ddot q+\frac{2}{C}\dot q
+Bl\left(\frac{3Bl}{2m}\dot q+\frac{F_0}{2m}\cos\omega t\right)=0.
\]

展开并整理：

\[
2r\ddot q+\left(\frac2C+\frac{3B^2l^2}{2m}\right)\dot q
=-\frac{BlF_0}{2m}\cos\omega t.
\]

为简化记号，定义

\[
\boxed{D=\frac2C+\frac{3B^2l^2}{2m}}.
\]

于是

\[
2r\ddot q+D\dot q=-\frac{BlF_0}{2m}\cos\omega t.
\]

对时间积分一次。由于 \(q(0)=0\)、\(\dot q(0)=0\)，积分常数为零，得到

\[
2r\dot q+Dq=-\frac{BlF_0}{2m\omega}\sin\omega t.
\]

再定义

\[
\boxed{E=\frac{BlF_0}{2m\omega}},
\]

则标准形式为

\[
\boxed{2r\dot q+Dq=-E\sin\omega t}.
\]

\textbf{（5）求电荷的简谐特解}

设特解为

\[q_p=P\sin\omega t+Q\cos\omega t.\]

则

\[\dot q_p=P\omega\cos\omega t-Q\omega\sin\omega t.\]

代入 \(2r\dot q+Dq=-E\sin\omega t\)：

\[
2r(P\omega\cos\omega t-Q\omega\sin\omega t)
+D(P\sin\omega t+Q\cos\omega t)
=-E\sin\omega t.
\]

比较 \(\sin\omega t\) 与 \(\cos\omega t\) 的系数：

\[
\begin{cases}
DP-2r\omega Q=-E,\\
2r\omega P+DQ=0.
\end{cases}
\]

由第二式

\[Q=-\frac{2r\omega}{D}P.\]

代入第一式：

\[
DP+\frac{4r^2\omega^2}{D}P=-E.
\]

两边乘 \(D\)：

\[
(D^2+4r^2\omega^2)P=-ED.
\]

因此

\[
\boxed{P=-\frac{ED}{D^2+4r^2\omega^2}},
\]

并有

\[
\boxed{Q=\frac{2r\omega E}{D^2+4r^2\omega^2}}.
\]

齐次方程 \(2r\dot q+Dq=0\) 的解为

\[q_h=K_0e^{-Dt/(2r)}.\]

所以总解为

\[
q(t)=P\sin\omega t+Q\cos\omega t+K_0e^{-Dt/(2r)}.
\]

由 \(q(0)=0\)：

\[0=Q+K_0,\]

故 \(K_0=-Q\)。最终

\[
\boxed{
q(t)=P\sin\omega t+Q\left(\cos\omega t-e^{-Dt/(2r)}\right)
},
\]

其中

\[
P=-\frac{ED}{D^2+4r^2\omega^2},
\qquad
Q=\frac{2r\omega E}{D^2+4r^2\omega^2}.
\]

直接求导：

\[
\boxed{
\dot q=P\omega\cos\omega t
+Q\left(\frac{D}{2r}e^{-Dt/(2r)}-\omega\sin\omega t\right)
}.
\]

把 \(t=0\) 代入可得

\[
\dot q(0)=P\omega+Q\frac{D}{2r}=0,
\]

与初始无电流相符。

\textbf{（6）回路电动势}

由

\[\mathcal E(t)=2r\dot q+\frac{2q}{C},\]

最简洁而完整的结果是

\[
\boxed{
\begin{aligned}
\mathcal E(t)={}&
2r\left[P\omega\cos\omega t
+Q\left(\frac{D}{2r}e^{-Dt/(2r)}-\omega\sin\omega t\right)\right]
\\
&+\frac2C\left[P\sin\omega t
+Q\left(\cos\omega t-e^{-Dt/(2r)}\right)\right]
\end{aligned}
}.
\]

若需要把各时间函数的系数完全展开，可写成

\[
\boxed{
\begin{aligned}
\mathcal E(t)
={}&2rP\omega\cos\omega t
+\left(\frac{2P}{C}-2rQ\omega\right)\sin\omega t\\
&+\frac{2Q}{C}\cos\omega t
+Q\left(D-\frac2C\right)e^{-Dt/(2r)}.
\end{aligned}}
\]

由于 \(D-2/C=3B^2l^2/(2m)\)，指数项系数还可继续化简。

\textbf{（7）两棒速度}

左棒速度已经得到：

\[
\boxed{
v_2(t)=-\frac{Bl}{m}\left[P\sin\omega t+Q\left(\cos\omega t-e^{-Dt/(2r)}\right)\right]
}.
\]

利用总动量关系

\[2mv_1+mv_2=\frac{F_0}{\omega}\sin\omega t,\]

解出 \(v_1\)：

\[
\begin{aligned}
v_1
&=\frac{F_0}{2m\omega}\sin\omega t-\frac12v_2\\
&=\frac{F_0}{2m\omega}\sin\omega t+\frac{Bl}{2m}q(t).
\end{aligned}
\]

故

\[
\boxed{
v_1(t)=\frac{F_0}{2m\omega}\sin\omega t
+\frac{Bl}{2m}\left[P\sin\omega t+Q\left(\cos\omega t-e^{-Dt/(2r)}\right)\right]
}.
\]

\textbf{量纲与极限核验}

\(D\) 的两项量纲必须一致。\(2/C\) 的量纲是伏特每库仑；\(B^2l^2/m\)
乘电荷后也给出电压，因此 \(Dq\) 与 \(2r\dot q\) 都是电压。

当 \(B\to0\) 时，\(E\to0\)，所以
\(P,Q,q\to0\)。此时左棒始终静止，右棒满足

\[2m\dot v_1=F_0\cos\omega t,\]

即

\[v_1=\frac{F_0}{2m\omega}\sin\omega t,\]

与通式一致。

当 \(t\gg2r/D\)
时，指数暂态消失，只剩与外力同频的稳态响应。电阻越大，衰减时间 \(2r/D\)
越长；但稳态相位与幅值也同时变化，不能只从衰减时间判断电流大小。

\newpage

\section{第 4 题}\label{ux7b2c-4-ux9898}

一个质量为 \(M\)、半径为 \(r\) 的中性孤立导体球固定在 \(x\)--\(y\)
平面的原点。一个质量为 \(m\ll M\)、电荷量为 \(q\) 的点电荷，从 \(x\)
轴负方向的远处以速度 \(v_0\) 沿直线 \(y=l\) 入射，其中
\(r\ll l\)。忽略感应电荷移动造成的能量损失，忽略相对论、电磁辐射以及万有引力。

\begin{enumerate}
\def\labelenumi{\arabic{enumi}.}
\tightlist
\item
  当点电荷到球心的距离为 \(h\) 时，求点电荷所受静电力大小 \(F(h)\)。
\item
  当 \(v_0\) 很大、轨迹只发生小角度偏转时，保留最低阶近似，求偏转角
  \(\alpha\)。
\item
  当 \(v_0\) 较小时，点电荷可能被导体球俘获。求临界入射速度 \(v_c\)。
\end{enumerate}

\begin{figure}[H]
\centering
\includegraphics[width=0.79\linewidth,height=\textheight,keepaspectratio,alt={图 4-1　点电荷掠过中性孤立导体球的散射与俘获。}]{../figures/question-04.pdf}
\caption*{图 4-1　点电荷掠过中性孤立导体球的散射与俘获。}
\end{figure}

\subsection{一、题型定位}\label{ux4e00ux9898ux578bux5b9aux4f4d-3}

中性孤立导体球的像电荷系统必须保持总感应电荷为零，因此接地球的单一负像电荷并不够。得到球外势能后，再求散射与俘获的临界条件。

\subsection{二、知识点}\label{ux4e8cux77e5ux8bc6ux70b9-3}

\begin{itemize}
\tightlist
\item
  中性孤立导体球的电像法。
\item
  接地球与中性孤立球像电荷配置的区别。
\item
  感应相互作用势能中的 \(1/2\) 因子。
\item
  高速小角散射的直线轨道近似。
\item
  标准积分 \(\int_{-\infty}^{\infty}\mathrm{d}x/(x^2+l^2)^3\)。
\item
  中心力角动量守恒。
\item
  有效势 \(U_{\rm eff}=L^2/(2m\rho^2)+U(\rho)\)。
\item
  临界俘获的双条件：\(E=U_{\rm eff}(\rho_0)\) 与
  \(U'_{\rm eff}(\rho_0)=0\)。
\end{itemize}

\subsection{三、物理图像}\label{ux4e09ux7269ux7406ux56feux50cf-3}

点电荷靠近中性导体球时，近侧出现异号感应电荷，远侧出现同号补偿分布。整体作用仍为吸引，但比接地导体球的吸引略弱。用像法表示时，在球内放置

\[q'=-\frac rhq\]

和

\[q_0=+\frac rhq.\]

前者位于球心到点电荷连线上、距球心 \(r^2/h\) 处，后者位于球心。

高速时，粒子来不及显著改变路径，吸引力只积累一个小横向冲量。低速时，粒子有足够时间被吸引；但若角动量较大，离心势会形成势垒。临界速度对应粒子刚好爬到势垒顶，径向速度在该处趋于零。

\subsection{四、解题入口}\label{ux56dbux89e3ux9898ux5165ux53e3-3}

第一问直接构造两个像电荷并分别计算对真实点电荷的力。第二问在未扰动轨迹上取

\[h=\sqrt{x^2+l^2},\qquad \mathrm{d}t=\frac{\mathrm{d}x}{v_0},\]

把径向吸引力投影到 \(y\)
方向后积分。第三问先由第一问积分得到真实相互作用势能，再写角动量

\[L_z=mv_cl\]

和有效势。临界条件必须同时使用势垒驻点与能量相等，不能只令径向加速度为零。

\subsection{五、解析}\label{ux4e94ux89e3ux6790-3}

图 4-2 左侧画出球内的负像电荷 \(-qr/h\)、球心补偿像电荷 \(+qr/h\)
及真实电荷所受合力；右侧画出角动量给出的离心项与感应吸引项共同形成的有效势垒。第（1）---（3）步使用左图，第（4）---（6）步使用右图。

\begin{figure}[H]
\centering
\includegraphics[width=0.9\linewidth,height=\textheight,keepaspectratio,alt={图 4-2　中性导体球的像电荷配置与临界俘获的有效势垒。}]{../figures/solution-04.pdf}
\caption*{图 4-2　中性导体球的像电荷配置与临界俘获的有效势垒。}
\end{figure}

\textbf{（1）中性孤立导体球的像电荷} 按图 4-2 左侧，设真实点电荷距球心为
\(h>r\)。负像电荷位于距球心 \(r^2/h\) 处；为了使半径为 \(r\)
的球面成为等势面，先写出接地球所需像电荷

\[
q'=-\frac rhq,
\]

其位置距球心

\[
x'=\frac{r^2}{h}.
\]

这个像电荷所代表的球面总感应电荷为
\(q'\)，并不满足中性条件。故再在球心放置补偿像电荷

\[
q_0=-q'=+\frac rhq.
\]

球心电荷在球面上产生常量电势，不破坏等势条件，同时使总像电荷为零。

真实电荷到负像电荷的距离为

\[
d=h-\frac{r^2}{h}=\frac{h^2-r^2}{h}.
\]

负像电荷对真实电荷的吸引力大小为

\[
\begin{aligned}
F_-
&=k_e\frac{|qq'|}{d^2}\\
&=k_e\frac{q^2(r/h)}{[(h^2-r^2)/h]^2}\\
&=k_eq^2\frac{rh}{(h^2-r^2)^2}.
\end{aligned}
\]

其中

\[k_e=\frac{1}{4\pi\varepsilon_0}.\]

球心正像电荷对真实电荷产生排斥力，大小为

\[
F_+=k_e\frac{q q_0}{h^2}
=k_eq^2\frac{r}{h^3}.
\]

净力仍指向球心，大小

\[
F(h)=F_--F_+.
\]

代入：

\[
F(h)=k_eq^2r\left[\frac{h}{(h^2-r^2)^2}-\frac1{h^3}\right].
\]

通分：

\[
\begin{aligned}
\frac{h}{(h^2-r^2)^2}-\frac1{h^3}
&=\frac{h^4-(h^2-r^2)^2}{h^3(h^2-r^2)^2}\\
&=\frac{h^4-(h^4-2h^2r^2+r^4)}{h^3(h^2-r^2)^2}\\
&=\frac{r^2(2h^2-r^2)}{h^3(h^2-r^2)^2}.
\end{aligned}
\]

所以

\[
\boxed{
F(h)=\frac{k_eq^2r^3(2h^2-r^2)}{h^3(h^2-r^2)^2}
},
\]

方向沿径向指向球心。

当 \(h\gg r\) 时，保留最低阶：

\[
F(h)\simeq\frac{2k_eq^2r^3}{h^5}.
\]

这正是点电荷与其在中性小导体球上诱导偶极矩相互作用的 \(h^{-5}\) 力律。

\textbf{（2）高速小角散射的横向冲量}

在小角近似下，粒子速度大小近似保持 \(v_0\)，轨迹可用未偏转直线 \(y=l\)
代替。粒子位于 \((x,l)\) 时

\[
h=\sqrt{x^2+l^2}.
\]

由于 \(r\ll l\)，整个主要作用区都有 \(h\gg r\)，故采用

\[F(h)\simeq\frac{2k_eq^2r^3}{h^5}.\]

径向力指向原点。它在垂直于入射方向的 \(y\) 分量大小为

\[
F_y=F(h)\frac{l}{h}.
\]

于是

\[
F_y\simeq\frac{2k_eq^2r^3l}{h^6}
=\frac{2k_eq^2r^3l}{(x^2+l^2)^3}.
\]

横向冲量大小

\[
I_y=\int_{-\infty}^{\infty}F_y\,\mathrm{d}t.
\]

由 \(\mathrm{d}t=\mathrm{d}x/v_0\)：

\[
I_y=\frac{2k_eq^2r^3l}{v_0}
\int_{-\infty}^{\infty}\frac{\mathrm{d}x}{(x^2+l^2)^3}.
\]

计算积分。令 \(x=l\tan\phi\)，则

\[\mathrm{d}x=l\sec^2\phi\,\mathrm{d}\phi,\]

\[x^2+l^2=l^2\sec^2\phi.\]

因此

\[
\begin{aligned}
\int_{-\infty}^{\infty}\frac{\mathrm{d}x}{(x^2+l^2)^3}
&=\frac1{l^5}\int_{-\pi/2}^{\pi/2}\cos^4\phi\,\mathrm{d}\phi\\
&=\frac1{l^5}\cdot\frac{3\pi}{8}\\
&=\frac{3\pi}{8l^5}.
\end{aligned}
\]

所以

\[
\begin{aligned}
I_y
&=\frac{2k_eq^2r^3l}{v_0}\cdot\frac{3\pi}{8l^5}\\
&=\frac{3\pi k_eq^2r^3}{4v_0l^4}.
\end{aligned}
\]

小角度下

\[\alpha\simeq\frac{|\Delta p_y|}{p_x}=\frac{I_y}{mv_0}.\]

故

\[
\boxed{
\alpha=\frac{3\pi k_eq^2r^3}{4mv_0^2l^4}
}.
\]

代入 \(k_e=1/(4\pi\varepsilon_0)\)：

\[
\boxed{
\alpha=\frac{3q^2r^3}{16\varepsilon_0mv_0^2l^4}
}.
\]

偏转方向指向导体球。

\textbf{（3）由静电力求相互作用势能}

定义 \(U(\infty)=0\)。因为径向力为

\[\bm F=-F(h)\hat{\bm r},\]

所以

\[\frac{\mathrm{d}U}{\mathrm{d}h}=F(h).\]

对第一问的精确力积分，可得

\[
\boxed{
U(h)=-\frac{k_eq^2r^3}{2h^2(h^2-r^2)}
}.
\]

可直接求导核验：

\[
\begin{aligned}
\frac{\mathrm{d}}{\mathrm{d}h}\left[-\frac{k_eq^2r^3}{2h^2(h^2-r^2)}\right]
&=\frac{k_eq^2r^3(2h^2-r^2)}{h^3(h^2-r^2)^2}\\
&=F(h).
\end{aligned}
\]

势能中的 \(1/2\)
是感应电荷随真实电荷逐渐建立的结果，不能把真实电荷与两个像电荷的普通库仑势能直接相加而漏掉该因子。

\textbf{（4）建立有效势}

转看图 4-2
右侧。中心力使角动量守恒，图中的势垒正是离心势与感应吸引势之和。临界入射速度记为
\(v_c\)，远处冲量参数为 \(l\)，故角动量

\[
L_z=mv_cl.
\]

径向运动的有效势为

\[
\begin{aligned}
U_{\rm eff}(h)
&=\frac{L_z^2}{2mh^2}+U(h)\\
&=\frac{mv_c^2l^2}{2h^2}
-\frac{k_eq^2r^3}{2h^2(h^2-r^2)}.
\end{aligned}
\]

第一项是离心势，第二项是感应吸引势。远处
\(U_{\rm eff}\to0^+\)；靠近球面时吸引势趋于
\(-\infty\)，两者之间存在一个势垒最大值。

\textbf{（5）势垒驻点条件}

令图 4-2 右侧虚线所示的势垒顶位置为 \(h_0\)。由

\[\left.\frac{\mathrm{d}U_{\rm eff}}{\mathrm{d}h}\right|_{h_0}=0\]

得到

\[
\frac{mv_c^2l^2}{h_0^3}
=\frac{k_eq^2r^3(2h_0^2-r^2)}{h_0^3(h_0^2-r^2)^2}.
\]

约去 \(h_0^{-3}\)：

\[
\boxed{
mv_c^2l^2
=k_eq^2r^3\frac{2h_0^2-r^2}{(h_0^2-r^2)^2}
}.
\]

\textbf{（6）临界能量条件}

远处总能量为

\[E=\frac12mv_c^2.\]

临界轨道刚好到达势垒顶，故

\[
\frac12mv_c^2
=\frac{mv_c^2l^2}{2h_0^2}
-\frac{k_eq^2r^3}{2h_0^2(h_0^2-r^2)}.
\]

两边乘 \(2h_0^2\)：

\[
mv_c^2h_0^2
=mv_c^2l^2-\frac{k_eq^2r^3}{h_0^2-r^2}.
\]

用驻点条件把 \(k_eq^2r^3\) 消去：

\[
k_eq^2r^3
=mv_c^2l^2\frac{(h_0^2-r^2)^2}{2h_0^2-r^2}.
\]

因此

\[
\frac{k_eq^2r^3}{h_0^2-r^2}
=mv_c^2l^2\frac{h_0^2-r^2}{2h_0^2-r^2}.
\]

代回能量式并约去 \(mv_c^2\)：

\[
h_0^2=l^2\left(1-\frac{h_0^2-r^2}{2h_0^2-r^2}\right).
\]

括号内通分：

\[
1-\frac{h_0^2-r^2}{2h_0^2-r^2}
=\frac{h_0^2}{2h_0^2-r^2}.
\]

所以

\[
h_0^2=l^2\frac{h_0^2}{2h_0^2-r^2}.
\]

约去 \(h_0^2\)：

\[
2h_0^2-r^2=l^2.
\]

故

\[
\boxed{h_0^2=\frac{l^2+r^2}{2}}.
\]

将 \(2h_0^2-r^2=l^2\) 和

\[h_0^2-r^2=\frac{l^2-r^2}{2}\]

代回驻点条件：

\[
mv_c^2l^2
=k_eq^2r^3\frac{l^2}{[(l^2-r^2)/2]^2}.
\]

右侧化简为

\[
k_eq^2r^3\frac{4l^2}{(l^2-r^2)^2}.
\]

约去 \(l^2\)：

\[
mv_c^2=\frac{4k_eq^2r^3}{(l^2-r^2)^2}.
\]

因此精确临界速度为

\[
\boxed{
v_c=\sqrt{\frac{4k_eq^2r^3}{m(l^2-r^2)^2}}
}.
\]

题设 \(r\ll l\)，故 \(l^2-r^2\simeq l^2\)，得到最低阶结果

\[
\boxed{
v_c\simeq\sqrt{\frac{4k_eq^2r^3}{ml^4}}
}.
\]

在冲量参数 \(l\) 固定时，角动量 \(L=mv_0l\) 也随速度增大。势垒高度随
\(v_0\) 的增长快于入射动能：当 \(v_0<v_c\)
时，入射能量高于该角动量对应的势垒，粒子能够越过势垒并落向导体球，因而被俘获；当
\(v_0>v_c\)
时，对应离心势垒高于入射能量，粒子在球外出现转向点并散射逃逸。

\textbf{与近似 \(h^{-5}\) 力的交叉核验}

若从一开始就采用

\[F\simeq\frac{2k_eq^2r^3}{h^5},\]

则

\[U\simeq-\frac{k_eq^2r^3}{2h^4}.\]

有效势

\[U_{\rm eff}\simeq\frac{mv_c^2l^2}{2h^2}-\frac{k_eq^2r^3}{2h^4}.\]

驻点条件给出

\[mv_c^2l^2h_0^2=2k_eq^2r^3,\]

能量条件给出 \(h_0^2=l^2/2\)。联立仍得到

\[v_c^2=\frac{4k_eq^2r^3}{ml^4}.\]

主阶结果与精确式一致；有限半径修正包含在 \(l^2-r^2\) 中。

\newpage

\section{第 5 题}\label{ux7b2c-5-ux9898}

水平面上固定一个半径为 \(R\)
的光滑圆筒。圆筒中心固定一根质量和原长均可忽略、劲度系数为 \(k\)
的轻弹簧，弹簧另一端连接质量为 \(m\)
的光滑质点。质点只能在圆筒内部运动。碰撞均为完全弹性碰撞。

\begin{enumerate}
\def\labelenumi{\arabic{enumi}.}
\item
  将质点置于 \(x_0=12R/13\)、\(y_0=0\) 处，并给它沿 \(+y\)
  方向的初速度。若要求质点始终不与圆筒碰撞，最大初速度写成

  \[v_{\max}=\alpha_1\sqrt{\frac{kR^2}{m}},\]

  求 \(\alpha_1\)。
\item
  若初速度大于
  \(v_{\max}\)，并要求碰撞后轨迹闭合，且各条轨迹``边''除公共``顶点''外不相交，所有可能轨迹中顶点数最多为多少？
\item
  若闭合轨迹有 \(4\) 个顶点，初速度写成

  \[v_4=\alpha_4\sqrt{\frac{kR^2}{m}},\]

  求 \(\alpha_4\)。
\item
  在第 3 问条件下，运动周期写成

  \[T=\beta_1\sqrt{\frac{m}{k}},\]

  求 \(\beta_1\)。
\end{enumerate}

\begin{figure}[H]
\centering
\includegraphics[width=0.45\linewidth,height=\textheight,keepaspectratio,alt={图 5-1　圆筒内由中心弹簧约束的二维质点运动。}]{../figures/question-05.pdf}
\caption*{图 5-1　圆筒内由中心弹簧约束的二维质点运动。}
\end{figure}

\subsection{一、题型定位}\label{ux4e00ux9898ux578bux5b9aux4f4d-4}

圆筒内的质点在相邻碰撞间作二维简谐运动，撞壁只改变径向速度。把每段椭圆轨迹与圆壁相交，就能追踪碰撞序列和顶点条件。

\subsection{二、知识点}\label{ux4e8cux77e5ux8bc6ux70b9-4}

\begin{itemize}
\tightlist
\item
  二维各向同性简谐运动的参数方程。
\item
  椭圆轨迹半轴与初始速度的关系。
\item
  圆形光滑壁的镜面弹性反射。
\item
  径向冲量不改变角动量。
\item
  闭合非自交轨迹与内接正多边形顶点。
\item
  由投影半轴得到顶点圆心角下限。
\item
  相位时间与完整轨道周期。
\item
  ``接触''\,``碰撞''与``切向掠过''的区分。
\end{itemize}

\subsection{三、物理图像}\label{ux4e09ux7269ux7406ux56feux50cf-4}

设

\[\omega_0=\sqrt{\frac{k}{m}}.\]

初始时 \(x=x_0=12R/13\)、\(\dot x=0\)，所以

\[x=x_0\cos\omega_0t.\]

初始时 \(y=0\)、\(\dot y=v_0\)，所以

\[y=\frac{v_0}{\omega_0}\sin\omega_0t.\]

若 \(v_0/\omega_0\le R\)，整条椭圆在圆内；若
\(v_0/\omega_0>R\)，椭圆会切穿圆壁，真实运动就在交点发生反射，并转入另一条全等椭圆弧。

对于四顶点轨迹，相邻碰撞点构成正方形。初始点位于右侧相邻两顶点之间那条椭圆弧的对称中点。若从初始点到第一个碰撞点经历相位
\(\tau\)，那么相邻两碰撞点之间经历 \(2\tau\)，完整四边轨迹经历
\(8\tau\)。

\subsection{四、解题入口}\label{ux56dbux89e3ux9898ux5165ux53e3-4}

先写无碰撞自由解

\[
x=a\cos\omega_0t,
\qquad
y=b\sin\omega_0t,
\]

其中

\[a=\frac{12}{13}R,\qquad b=\frac{v_0}{\omega_0}.\]

第一问只比较 \(a,b,R\)。第二问取初始点两侧的相邻碰撞顶点，它们关于 \(x\)
轴对称，若每个顶点与 \(x\) 轴夹角为 \(\varphi\)，则
\(R\cos\varphi\le a\)。第三问把四顶点坐标 \((R/\sqrt2,R/\sqrt2)\)
代入椭圆方程。第四问由 \(x=a\cos\omega_0t\) 求首次碰撞相位。

\subsection{五、解析}\label{ux4e94ux89e3ux6790-4}

图 5-2
左侧把四个碰撞点画成圆内接正方形，并用蓝色画出碰撞之间的椭圆弧；右侧用同一相位变量画出
\(x\)、\(y\)
的简谐变化。第（3）---（4）步使用左图确定顶点几何，第（5）步使用右图计算相位间隔和周期。

\begin{figure}[H]
\centering
\includegraphics[width=0.9\linewidth,height=\textheight,keepaspectratio,alt={图 5-2　四顶点闭合轨迹的椭圆弧、碰撞点与相位关系。}]{../figures/solution-05.pdf}
\caption*{图 5-2　四顶点闭合轨迹的椭圆弧、碰撞点与相位关系。}
\end{figure}

\textbf{（1）无碰撞轨迹的最大初速度} 质点在碰撞前的运动方程为

\[
\ddot x+\omega_0^2x=0,
\qquad
\ddot y+\omega_0^2y=0,
\]

其中

\[\omega_0=\sqrt{\frac{k}{m}}.\]

由初始条件

\[x(0)=\frac{12}{13}R,\quad \dot x(0)=0\]

得到

\[
x(t)=\frac{12}{13}R\cos\omega_0t.
\]

由

\[y(0)=0,\quad \dot y(0)=v_0\]

得到

\[
y(t)=\frac{v_0}{\omega_0}\sin\omega_0t.
\]

消去时间：

\[
\boxed{
\frac{x^2}{(12R/13)^2}+
\frac{y^2}{(v_0/\omega_0)^2}=1
}.
\]

这是以圆心为中心、半轴分别为

\[a=\frac{12}{13}R,\qquad b=\frac{v_0}{\omega_0}\]

的椭圆。

由于 \(a<R\)，是否碰壁只取决于另一半轴 \(b\)。要求整条椭圆位于半径为
\(R\) 的圆内，必须且只需

\[b\le R.\]

即

\[
\frac{v_0}{\omega_0}\le R.
\]

所以

\[
v_0\le\omega_0R
=R\sqrt{\frac{k}{m}}
=\sqrt{\frac{kR^2}{m}}.
\]

因此最大速度

\[
v_{\max}=\sqrt{\frac{kR^2}{m}},
\]

故

\[
\boxed{\alpha_1=1}.
\]

当取等号时，椭圆在 \(y=\pm R\)
处与圆壁相切，相对法向速度为零。理想模型中这是临界接触而不是具有有限冲量的碰撞。

\textbf{（2）碰撞前后保持的物理量}

圆壁光滑，碰撞冲量沿碰撞点的半径方向。设碰撞前速度分解为

\[\bm v^- =v_r^-\hat{\bm r}+v_t^-\hat{\bm t}.\]

完全弹性镜面反射给出

\[
v_r^+=-v_r^-,\qquad v_t^+=v_t^-.
\]

所以速度大小保持：

\[|\bm v^+|=|\bm v^-|.\]

碰撞位置不变，弹簧势能也不变，因此总机械能守恒。又因为冲量沿半径，对圆心的冲量矩为零，所以角动量

\[L_z=mRv_t\]

也守恒。

对各向同性谐振子，能量决定半轴平方和，角动量决定半轴乘积。因此碰撞后的下一段轨迹仍是一条半轴长度相同的椭圆，只是椭圆主轴方向改变。

\textbf{（3）相邻顶点圆心角的下限}

按图 5-2 左侧，考虑包含初始点 \((a,0)\) 的第一段蓝色椭圆弧，其中

\[a=\frac{12}{13}R.\]

这条弧关于 \(x\) 轴对称，初始点是该弧在 \(x\)
方向的最右端。设它前后相邻的两个碰撞点关于 \(x\) 轴对称，极角分别为
\(+\varphi\) 与 \(-\varphi\)。每个碰撞点的横坐标为

\[x_v=R\cos\varphi.\]

另一方面，椭圆上任意点满足 \(|x|\le a\)，因此

\[
R\cos\varphi\le a=\frac{12}{13}R.
\]

约去 \(R\)：

\[
\cos\varphi\le\frac{12}{13}.
\]

在 \(0<\varphi<\pi/2\) 内，余弦函数单调递减，所以

\[
\boxed{\varphi\ge\arccos\frac{12}{13}}.
\]

相邻两顶点的圆心角为

\[
\Delta\Phi=2\varphi,
\]

故

\[
\boxed{\Delta\Phi\ge2\arccos\frac{12}{13}}.
\]

若闭合轨迹有 \(N\)
个顶点，且各边除公共顶点外不相交，则顶点依次绕圆周一周，满足

\[N\Delta\Phi=2\pi.\]

于是

\[
N\le\frac{2\pi}{2\arccos(12/13)}
=\frac{\pi}{\arccos(12/13)}.
\]

数值上

\[
\arccos\frac{12}{13}\simeq0.39479\ \mathrm{rad}
=22.62^\circ,
\]

所以

\[
\frac{\pi}{\arccos(12/13)}\simeq7.96.
\]

\(N\) 必须为整数，因此

\[
\boxed{N_{\max}=7}.
\]

七边形所需半角为 \(\pi/7\simeq25.71^\circ\)，大于
\(22.62^\circ\)，可以通过选择适当的 \(b\) 实现；八边形所需半角为
\(22.5^\circ\)，小于允许下限，不能实现。

\textbf{（4）四顶点闭合轨迹的初速度}

图 5-2 左侧四个顶点组成圆内接正方形。取第一象限碰撞点，其坐标为

\[
x_v=y_v=\frac{R}{\sqrt2}.
\]

初始椭圆方程是

\[
\frac{x^2}{a^2}+\frac{y^2}{b^2}=1,
\]

其中

\[a=\frac{12}{13}R,\qquad b=\frac{v_4}{\omega_0}.\]

将碰撞点代入：

\[
\frac{R^2/2}{(12R/13)^2}
+\frac{R^2/2}{b^2}=1.
\]

第一项为

\[
\frac{R^2/2}{144R^2/169}=\frac{169}{288}.
\]

故

\[
\frac{R^2}{2b^2}=1-\frac{169}{288}=\frac{119}{288}.
\]

两边乘以 \(2b^2\)：

\[
R^2=\frac{119}{144}b^2.
\]

因此

\[
b^2=\frac{144}{119}R^2,
\]

\[
\boxed{b=\frac{12}{\sqrt{119}}R}.
\]

由 \(b=v_4/\omega_0\)：

\[
v_4=\omega_0b
=\sqrt{\frac{k}{m}}\frac{12R}{\sqrt{119}}.
\]

整理为题设形式：

\[
\boxed{
v_4=\frac{12}{\sqrt{119}}\sqrt{\frac{kR^2}{m}}
}.
\]

故

\[
\boxed{\alpha_4=\frac{12}{\sqrt{119}}\simeq1.100}.
\]

它确实大于 \(\alpha_1=1\)，符合必须发生碰撞的要求。

\textbf{（5）四顶点轨迹的周期}

转看图 5-2 右侧：第一次碰撞相位记为
\(\tau\)，相邻碰撞点的相位间隔由波形对称性为 \(2\tau\)。从初始点
\((a,0)\) 到第一象限碰撞点，先由

\[x(t)=a\cos\omega_0t\]

和

\[x(\tau)=\frac{R}{\sqrt2}\]

得到

\[
a\cos\omega_0\tau=\frac{R}{\sqrt2}.
\]

代入 \(a=12R/13\)：

\[
\frac{12R}{13}\cos\omega_0\tau=\frac{R}{\sqrt2}.
\]

约去 \(R\)：

\[
\cos\omega_0\tau=\frac{13}{12\sqrt2}.
\]

因此

\[
\boxed{
\tau=\frac1{\omega_0}\arccos\frac{13}{12\sqrt2}
}.
\]

初始点是右侧两相邻碰撞顶点之间椭圆弧的对称中点。因此，从一个碰撞顶点到下一个碰撞顶点需时
\(2\tau\)。四条边组成完整闭合轨迹，所以

\[
T=4(2\tau)=8\tau.
\]

于是

\[
\begin{aligned}
T
&=\frac8{\omega_0}\arccos\frac{13}{12\sqrt2}\\
&=8\arccos\frac{13}{12\sqrt2}\sqrt{\frac{m}{k}}.
\end{aligned}
\]

故

\[
\boxed{
\beta_1=8\arccos\frac{13}{12\sqrt2}
\simeq5.585
}.
\]

\textbf{碰撞反射为何仍给出全等椭圆}

二维各向同性谐振子的能量可写为

\[E=\frac12m\omega_0^2(a^2+b^2),\]

角动量大小为

\[|L_z|=m\omega_0ab.\]

碰撞保持 \(E\) 与 \(|L_z|\)，因此 \(a^2+b^2\) 与 \(ab\)
都不变。\(a^2,b^2\) 是方程

\[z^2-(a^2+b^2)z+a^2b^2=0\]

的两根，所以无序半轴组 \(\{a,b\}\)
不变。碰撞只改变椭圆主轴方向，不改变形状。

\textbf{周期系数数值核验}

\[
\frac{13}{12\sqrt2}\simeq0.76603,
\]

\[
\arccos(0.76603)\simeq0.69815,
\]

因此

\[\beta_1\simeq8\times0.69815\simeq5.585.\]

\newpage

\section{第 6 题}\label{ux7b2c-6-ux9898}

以下各问均在非相对论框架内讨论，并取

\[\frac{1}{\varepsilon_0\mu_0}=c^2.\]

\begin{enumerate}
\def\labelenumi{\arabic{enumi}.}
\item
  基础场分布：

  \begin{enumerate}
  \def\labelenumii{\arabic{enumii}.}
  \tightlist
  \item
    真空中有一无限大均匀带电平面，面电荷密度为 \(\sigma\)，平面为
    \(x\)--\(y\) 平面。求电场分布。
  \item
    真空中有一沿 \(z\) 轴的无限长均匀带电直线，线电荷密度为
    \(\lambda\)。求距导线 \(r\) 处的电场强度。
  \item
    上述无限大带电平面以速度 \(\bm v=v\hat{\bm x}\)
    运动。求电场和磁场分布，并用 \(\bm E\) 表示 \(\bm B\)。
  \end{enumerate}
\item
  静止参考系中存在

  \[\bm E=E_x\hat{\bm x}+E_y\hat{\bm y}+E_z\hat{\bm z},\]

  且 \(\bm B=0\)。另一参考系以速度 \(\bm v=v\hat{\bm z}\)
  运动。求低速近似下该系中的 \(\bm E'\) 与 \(\bm B'\)。
\item
  沿 \(z\) 轴放置线电荷密度为 \(+\lambda\)
  的静止正电荷，同时线电荷密度为 \(-\lambda\) 的负电荷以速度
  \(\bm v=v\hat{\bm z}\) 运动。

  \begin{enumerate}
  \def\labelenumii{\arabic{enumii}.}
  \tightlist
  \item
    求静止观察者看到的距轴线 \(r\) 处的 \(\bm E\) 与 \(\bm B\)。
  \item
    另一观察者以 \(\bm v=v\hat{\bm z}\)
    运动。若仍按经典速度叠加，求其看到的 \(\bm E'\) 与 \(\bm B'\)。
  \item
    有一电荷量为 \(q\)
    的质点与运动观察者共同运动。为了使两个参考系计算出的质点受力相同，设只在运动参考系中额外出现线电荷密度
    \(\Delta\lambda\)。求 \(\Delta\lambda/\lambda\)。
  \end{enumerate}
\end{enumerate}

\begin{figure}[H]
\centering
\includegraphics[width=0.9\linewidth,height=\textheight,keepaspectratio,alt={图 6-1　带电平面、带电直线与运动参考系中的电磁场。}]{../figures/question-06.pdf}
\caption*{图 6-1　带电平面、带电直线与运动参考系中的电磁场。}
\end{figure}

\subsection{一、题型定位}\label{ux4e00ux9898ux578bux5b9aux4f4d-5}

先由高斯定律和安培定律求静止源的电磁场，再比较不同惯性系对同一粒子的受力。低速展开中出现的二阶差异指向经典速度叠加的局限。

\subsection{二、知识点}\label{ux4e8cux77e5ux8bc6ux70b9-5}

\begin{itemize}
\tightlist
\item
  高斯定律求无限平面与无限直线电场。
\item
  面电流密度 \(\bm K=\sigma\bm v\)。
\item
  线电流 \(I=\lambda v\)，负电荷运动时的常规电流方向。
\item
  无限电流直线磁场 \(B=\mu_0I/(2\pi r)\)。
\item
  低速场变换
  \(\bm E'\simeq\bm E\)、\(\bm B'\simeq\bm B-\bm v\times\bm E/c^2\)。
\item
  洛伦兹力 \(q(\bm E+\bm u\times\bm B)\)。
\item
  \(\varepsilon_0\mu_0=1/c^2\)。
\item
  一阶速度场与二阶密度修正的层级关系。
\end{itemize}

\subsection{三、物理图像}\label{ux4e09ux7269ux7406ux56feux50cf-5}

无限带电平面的电场垂直平面并在两侧反向。平面沿自身方向运动后形成面电流，磁场平行于平面且在两侧反向。无限带电直线的电场沿径向，运动线电荷形成沿轴电流，磁场沿环向。

第三问中，静止系看到总电荷密度为零，所以没有电场；但负电荷运动形成电流，所以有磁场。与负电荷共同运动的参考系仍按经典叠加得到同样大小的电流，只是改由正电荷反向运动产生。共同运动的试探粒子在该系静止，磁力消失。要让两系受力一致，只能增加一个径向电场，这等价于出现微小的净线电荷密度。

\subsection{四、解题入口}\label{ux56dbux89e3ux9898ux5165ux53e3-5}

第一部分全部用对称高斯面或安培环路。带电平面的磁场也可先写

\[\bm K=\sigma v\hat{\bm x},\]

再用表面电流边界条件求两侧磁场。

第二部分直接从低速变换

\[
\bm E'=\bm E,
\qquad
\bm B'=-\frac{\bm v\times\bm E}{c^2}
\]

展开全部分量。

第三部分先明确常规电流方向。线电荷 \(-\lambda\) 以 \(+v\hat z\)
运动，对应

\[I=-\lambda v.\]

在圆柱坐标中取 \(\hat{\bm r},\hat{\bm\varphi},\hat{\bm z}\)
右手系，再比较试探粒子在两系的径向力。

\subsection{五、解析}\label{ux4e94ux89e3ux6790-5}

图 6-2
左侧画出静止系中的负电流、环向磁场和随负电荷同速运动的试探电荷；右侧画出共同运动系中试探电荷静止后，为补偿磁力而必须出现的径向电场
\(E_{\Delta\lambda}\)。前两问先独立推导基本场公式，第（3a）---（3c）步再逐项使用这两幅图比较同一粒子的受力。

\begin{figure}[H]
\centering
\includegraphics[width=0.9\linewidth,height=\textheight,keepaspectratio,alt={图 6-2　两个参考系中的线电流、磁力与附加线电荷补偿。}]{../figures/solution-06.pdf}
\caption*{图 6-2　两个参考系中的线电流、磁力与附加线电荷补偿。}
\end{figure}

\textbf{（1a）无限大带电平面的电场} 平面为 \(z=0\)，面电荷密度为
\(\sigma\)。由平面对称性，电场只能垂直平面，并且两侧大小相等、方向相反。

取横截面积为
\(S\)、跨过平面的薄柱形高斯面。侧面无电通量，上下两底面的总通量为

\[\Phi_E=ES+ES=2ES.\]

高斯面所包围电荷为

\[Q_{\rm in}=\sigma S.\]

由高斯定律

\[2ES=\frac{\sigma S}{\varepsilon_0}.\]

约去 \(S\)：

\[E=\frac{\sigma}{2\varepsilon_0}.\]

若 \(\sigma>0\)，电场从平面向外，故完整矢量式为

\[
\boxed{
\bm E(z)=\frac{\sigma}{2\varepsilon_0}\operatorname{sgn}(z)\hat{\bm z}
}.
\]

\textbf{（1b）无限长带电直线的电场}

取半径为 \(r\)、长度为 \(L\)、与 \(z\)
轴同轴的圆柱高斯面。电场沿径向，大小在柱侧面上恒定。通量

\[\Phi_E=E(2\pi rL).\]

包围电荷

\[Q_{\rm in}=\lambda L.\]

高斯定律给出

\[E(2\pi rL)=\frac{\lambda L}{\varepsilon_0}.\]

约去 \(L\)：

\[
\boxed{
\bm E(r)=\frac{\lambda}{2\pi\varepsilon_0r}\hat{\bm r}
}.
\]

\(\lambda<0\) 时公式自动给出向内方向。

\textbf{（1c）运动带电平面的电场与磁场}

在非相对论近似中忽略面电荷密度的长度收缩修正，因此电场仍为

\[
\boxed{
\bm E(z)=\frac{\sigma}{2\varepsilon_0}\operatorname{sgn}(z)\hat{\bm z}
}.
\]

平面以 \(\bm v=v\hat{\bm x}\) 运动，形成面电流密度

\[
\bm K=\sigma v\hat{\bm x}.
\]

无限面电流两侧磁场大小均为

\[B=\frac{\mu_0K}{2}=\frac{\mu_0\sigma v}{2}.\]

用右手定则判断：在 \(z>0\)
一侧，\(\hat x\times\hat z=-\hat y\)，故磁场沿 \(-\hat y\)；在 \(z<0\)
一侧方向相反。因此

\[
\boxed{
\bm B(z)=-\frac{\mu_0\sigma v}{2}\operatorname{sgn}(z)\hat{\bm y}
}.
\]

又因为

\[
\frac1{c^2}=\varepsilon_0\mu_0,
\]

并且

\[
\bm v\times\bm E
=v\hat x\times\left[\frac{\sigma}{2\varepsilon_0}\operatorname{sgn}(z)\hat z\right]
=-\frac{\sigma v}{2\varepsilon_0}\operatorname{sgn}(z)\hat y,
\]

所以

\[
\boxed{
\bm B=\frac{\bm v\times\bm E}{c^2}
}.
\]

这个关系适用于此处``场源整体运动、在实验室系求场''的方向约定。

\textbf{（2）低速电磁场变换}

静止系中 \(\bm B=0\)，运动参考系速度为

\[\bm v=v\hat{\bm z}.\]

在 \(v\ll c\) 的最低阶近似下，

\[
\boxed{\bm E'\simeq\bm E},
\]

\[
\boxed{\bm B'\simeq-\frac{\bm v\times\bm E}{c^2}}.
\]

先计算叉积：

\[
\begin{aligned}
\bm v\times\bm E
&=v\hat z\times(E_x\hat x+E_y\hat y+E_z\hat z)\\
&=vE_x(\hat z\times\hat x)+vE_y(\hat z\times\hat y)\\
&=vE_x\hat y-vE_y\hat x.
\end{aligned}
\]

因此

\[
\begin{aligned}
\bm B'
&=-\frac{vE_x\hat y-vE_y\hat x}{c^2}\\
&=\frac{vE_y}{c^2}\hat x-\frac{vE_x}{c^2}\hat y.
\end{aligned}
\]

完整分量为

\[
\boxed{
\begin{aligned}
E_x'&=E_x,& E_y'&=E_y,& E_z'&=E_z,\\
B_x'&=\frac{vE_y}{c^2},&
B_y'&=-\frac{vE_x}{c^2},&
B_z'&=0.
\end{aligned}}
\]

由于题设电场含一般分量
\(E_x,E_y,E_z\)，答案必须补齐叉积产生的全部非零磁场分量。

\textbf{（3a）静止观察者看到的场}

按图 6-2 左侧，正线电荷 \(+\lambda\) 静止，负线电荷 \(-\lambda\) 沿
\(+z\) 运动并形成沿 \(-z\) 的常规电流。两者电荷密度代数和为

\[\lambda+(-\lambda)=0.\]

所以电场完全抵消：

\[
\boxed{\bm E=0}.
\]

负线电荷以 \(+v\hat z\) 运动。常规电流为电荷密度乘速度：

\[
I=(-\lambda)v=-\lambda v.
\]

负号说明常规电流沿 \(-\hat z\)。无限直线电流的磁场为

\[
\bm B=\frac{\mu_0I}{2\pi r}\hat{\bm\varphi}.
\]

因此

\[
\boxed{
\bm B=-\frac{\mu_0\lambda v}{2\pi r}\hat{\bm\varphi}
}.
\]

\textbf{（3b）与负电荷共同运动的观察者}

转到图 6-2 右侧。运动观察者速度为
\(+v\hat z\)，负电荷和试探电荷在该系静止；按经典速度叠加，正电荷则以速度

\[-v\hat z\]

运动。若仍忽略任何长度收缩，正、负线电荷密度大小仍相等，所以

\[
\boxed{\bm E'=0}.
\]

正线电荷产生的电流为

\[I'=\lambda(-v)=-\lambda v.\]

故

\[
\boxed{
\bm B'=-\frac{\mu_0\lambda v}{2\pi r}\hat{\bm\varphi}
}.
\]

经典处理给出了两系相同的磁场，却会在试探粒子受力上产生矛盾，因为粒子在一个系中运动、在另一个系中静止。

\textbf{（3c）比较同一试探粒子的受力}

图 6-2
的两幅图描述的是同一个试探粒子。它与运动观察者共同运动，所以在静止系中其速度为

\[\bm u=v\hat z.\]

静止系中 \(\bm E=0\)，粒子只受磁力：

\[
\bm F=q\bm u\times\bm B.
\]

代入

\[\bm B=-\frac{\mu_0\lambda v}{2\pi r}\hat\varphi.\]

利用柱坐标右手关系

\[\hat z\times\hat\varphi=-\hat r,\]

可得

\[
\hat z\times(-\hat\varphi)=+\hat r.
\]

所以对 \(q>0\)，磁力沿径向向外，大小

\[
\boxed{
F=\frac{q\mu_0\lambda v^2}{2\pi r}
}.
\]

在运动系中粒子静止，磁力为零。若额外出现线电荷密度
\(\Delta\lambda\)，它产生径向电场

\[
E_{\Delta}=\frac{\Delta\lambda}{2\pi\varepsilon_0r}.
\]

粒子所受电力大小为

\[
F'=qE_{\Delta}
=\frac{q\Delta\lambda}{2\pi\varepsilon_0r}.
\]

要求两个参考系给出相同的力：

\[
\frac{q\Delta\lambda}{2\pi\varepsilon_0r}
=\frac{q\mu_0\lambda v^2}{2\pi r}.
\]

约去 \(q/(2\pi r)\)：

\[
\frac{\Delta\lambda}{\varepsilon_0}=\mu_0\lambda v^2.
\]

因此

\[
\frac{\Delta\lambda}{\lambda}
=\varepsilon_0\mu_0v^2.
\]

利用 \(\varepsilon_0\mu_0=1/c^2\)：

\[
\boxed{
\frac{\Delta\lambda}{\lambda}=\frac{v^2}{c^2}
}.
\]

\(\Delta\lambda>0\)，因为运动系中需要一个向外电场来补偿静止粒子所缺少的向外磁力。

\textbf{近似阶数说明}

由运动电荷产生的磁场相对电场通常是一阶小量
\(B\sim(v/c^2)E\)。试探粒子的磁力还要再乘一个速度
\(v\)，因此磁力与对应电力之比达到
\(v^2/c^2\)。这正是第三问中净电荷密度修正的阶数。

完整狭义相对论会给出线密度的洛伦兹因子变化。低速展开

\[\gamma=\frac1{\sqrt{1-v^2/c^2}}\simeq1+\frac12\frac{v^2}{c^2}\]

后，正负电荷的密度差确实首先出现在二阶。题目用经典力矛盾反推出同一数量级，但不应把这里的简化推导误认为完整的相对论证明。

\newpage

\section{第 7 题}\label{ux7b2c-7-ux9898}

低压氢灯产生的可见光谱中有一条波长为 \(486.1\,\mathrm{nm}\)
的蓝绿色谱线。试求：

\begin{enumerate}
\def\labelenumi{\arabic{enumi}.}
\tightlist
\item
  一个该谱线光子的能量，用电子伏特表示。
\item
  放出该光子时，氢原子初态、末态分别处在第 \(i\) 激发态和第 \(f\)
  激发态，求 \(i,f\)。基态记为第 \(0\) 激发态。
\item
  电子从初态跃迁到末态的过程约持续
  \(1.5\times10^{-8}\,\mathrm{s}\)。根据能量---时间不确定关系，估算谱线的最小频宽。
\item
  为估计热运动造成的多普勒展宽，设一群氢原子只沿光线进入分光计的方向作一维运动，速率最大为
  \(700\,\mathrm{m/s}\)。求观察到的该蓝绿色谱线总频宽。
\item
  静止氢原子的谱线频率为
  \(\nu_0\)。若氢原子沿观测方向的速度满足一维麦克斯韦分布，热力学温度为
  \(T\)，求频偏 \(\Delta\nu\) 的方均根值 \(\Delta\nu_{\mathrm{rms}}\)。
\end{enumerate}

\begin{figure}[H]
\centering
\includegraphics[width=0.9\linewidth,height=\textheight,keepaspectratio,alt={图 7-1　低压氢灯可见光谱中的 486.1\textbackslash,\textbackslash mathrm\{nm\} 巴耳末线。}]{../figures/question-07.pdf}
\caption*{图 7-1　低压氢灯可见光谱中的 \(486.1\,\mathrm{nm}\) 巴耳末线。}
\end{figure}

\subsection{一、题型定位}\label{ux4e00ux9898ux578bux5b9aux4f4d-6}

486.1 nm
谱线先确定氢原子的跃迁能量，有限寿命给出自然线宽，热运动则带来多普勒展宽。几种宽度来源不同，需分别计算后再合成。

\subsection{二、知识点}\label{ux4e8cux77e5ux8bc6ux70b9-6}

\begin{itemize}
\tightlist
\item
  光子能量 \(E=h\nu=hc/\lambda\)。
\item
  氢原子能级 \(E_n=-13.6\,\mathrm{eV}/n^2\)。
\item
  激发态编号与主量子数的换算：第 \(i\) 激发态对应 \(n=i+1\)。
\item
  巴耳末系 \(n\to2\)。
\item
  能量---时间不确定关系 \(\Delta E\Delta t\gtrsim\hbar/2\)。
\item
  \(\Delta E=h\Delta\nu\)。
\item
  纵向多普勒频移。
\item
  一维麦克斯韦分布与 \(v_{x,\rm rms}=\sqrt{k_BT/m_H}\)。
\item
  总频宽、单侧频移与方均根宽度的区别。
\end{itemize}

\subsection{三、物理图像}\label{ux4e09ux7269ux7406ux56feux50cf-6}

\(486.1\,\mathrm{nm}\) 是氢原子巴耳末系的 \(H_\beta\) 线，对应电子从
\(n=4\) 跃迁到 \(n=2\)。光子能量等于两能级能量差。

若发光原子朝观察者运动，波峰到达得更密，观测频率升高；若背离观察者，频率降低。因此具有正负速度的原子群把单一频率展宽为一段频率区间。若速度服从高斯型的一维麦克斯韦分布，频偏也服从高斯分布，其标准差由速度标准差线性映射得到。

\subsection{四、解题入口}\label{ux56dbux89e3ux9898ux5165ux53e3-6}

第一问直接用常用常数

\[hc\simeq1240\,\mathrm{eV\,nm}.\]

第二问比较

\[13.6\left(\frac1{2^2}-\frac1{n^2}\right)\mathrm{eV}\]

与第一问的光子能量。第三问由 \(h\Delta\nu=\Delta E\)
换算自然线宽。第四问的速度区间是 \(-700\) 至
\(+700\,\mathrm{m/s}\)，总频宽约为单侧频移的两倍。第五问用

\[\Delta\nu\simeq-\frac{v_x}{c}\nu_0\]

平方取平均。

\subsection{五、解析}\label{ux4e94ux89e3ux6790-6}

图 7-2 左侧用能级图固定跃迁为
\(n=4\to2\)，右侧画出热运动速度分布线性映射后的频偏分布，并同时标出限速模型的总宽度
\(2\nu_0v_0/c\)。第（2）步使用左图，第（4）---（5）步使用右图；自然线宽则在第（3）步单独由寿命推导。

\begin{figure}[H]
\centering
\includegraphics[width=0.9\linewidth,height=\textheight,keepaspectratio,alt={图 7-2　氢原子 n=4\textbackslash to2 跃迁、自然线宽与热运动频偏分布。}]{../figures/solution-07.pdf}
\caption*{图 7-2　氢原子 \(n=4\to2\) 跃迁、自然线宽与热运动频偏分布。}
\end{figure}

\textbf{（1）光子能量} 光子能量

\[E_\gamma=h\nu_0=\frac{hc}{\lambda_0}.\]

代入

\[hc\simeq1240\,\mathrm{eV\,nm},\]

\[\lambda_0=486.1\,\mathrm{nm},\]

得到

\[
\begin{aligned}
E_\gamma
&=\frac{1240\,\mathrm{eV\,nm}}{486.1\,\mathrm{nm}}\\
&\simeq2.551\,\mathrm{eV}.
\end{aligned}
\]

所以

\[
\boxed{E_\gamma\simeq2.55\,\mathrm{eV}}.
\]

若用焦耳表示，

\[
E_\gamma\simeq2.55\times1.602\times10^{-19}\,\mathrm{J}
\simeq4.09\times10^{-19}\,\mathrm{J}.
\]

\textbf{（2）确定初末能级}

按图 7-2 左侧红色箭头，辐射从较高能级跃迁到 \(n=2\)。氢原子主量子数为
\(n\) 的能级是

\[
E_n=-\frac{13.6}{n^2}\,\mathrm{eV}.
\]

可见光巴耳末系的末态为 \(n_f=2\)。设初态为 \(n_i\)，则发射光子能量

\[
E_\gamma=13.6\left(\frac1{2^2}-\frac1{n_i^2}\right)\mathrm{eV}.
\]

尝试 \(n_i=4\)：

\[
\begin{aligned}
E_4-E_2
&=13.6\left(\frac14-\frac1{16}\right)\mathrm{eV}\\
&=13.6\cdot\frac3{16}\,\mathrm{eV}\\
&=2.55\,\mathrm{eV}.
\end{aligned}
\]

与第一问完全相符。因此跃迁为

\[
\boxed{n_i=4\to n_f=2}.
\]

题目把基态记为第 \(0\) 激发态，所以第 \(i\) 激发态对应主量子数

\[n=i+1.\]

于是

\[i=n_i-1=3,\]

\[f=n_f-1=1.\]

故

\[
\boxed{i=3,\qquad f=1}.
\]

\textbf{（3）跃迁有限寿命导致的自然线宽}

自然线宽来自激发态有限寿命，与右图的热运动多普勒展宽不是同一机制。采用能量---时间不确定关系的估算形式

\[
\Delta E\,\Delta t\simeq\frac{\hbar}{2}.
\]

由于

\[
\Delta E=h\Delta\nu,
\]

所以

\[
h\Delta\nu\,\Delta t\simeq\frac{\hbar}{2}.
\]

利用

\[\hbar=\frac{h}{2\pi},\]

得到

\[
h\Delta\nu\,\Delta t\simeq\frac{h}{4\pi}.
\]

约去 \(h\)：

\[
\Delta\nu\simeq\frac{1}{4\pi\Delta t}.
\]

代入

\[\Delta t=1.5\times10^{-8}\,\mathrm{s}:\]

\[
\begin{aligned}
\Delta\nu
&\simeq\frac{1}{4\pi\times1.5\times10^{-8}}\,\mathrm{Hz}\\
&\simeq5.31\times10^6\,\mathrm{Hz}.
\end{aligned}
\]

因此

\[
\boxed{\Delta\nu_{\rm nat}\simeq5.3\,\mathrm{MHz}}.
\]

这里的数值系数取决于对``线宽''与不确定关系等号的具体约定；竞赛估算中应明确采用了
\(\Delta E\Delta t\simeq\hbar/2\)。

\textbf{（4）最大速率模型下的多普勒总频宽}

静止原子的频率为

\[
\nu_0=\frac{c}{\lambda_0}.
\]

代入
\(c=2.998\times10^8\,\mathrm{m/s}\)、\(\lambda_0=486.1\times10^{-9}\,\mathrm m\)：

\[
\nu_0\simeq6.17\times10^{14}\,\mathrm{Hz}.
\]

设原子朝观察者运动时视线速度为 \(+v\)。纵向相对论多普勒公式为

\[
\nu_+=\nu_0\sqrt{\frac{1+\beta}{1-\beta}},
\qquad
\beta=\frac vc.
\]

背离观察者时

\[
\nu_-=\nu_0\sqrt{\frac{1-\beta}{1+\beta}}.
\]

因此两端之间的总频宽为

\[
\Delta\nu_{\rm D}=\nu_+-\nu_-.
\]

由于 \(v=700\,\mathrm{m/s}\ll c\)，有

\[
\sqrt{\frac{1+\beta}{1-\beta}}\simeq1+\beta,
\]

\[
\sqrt{\frac{1-\beta}{1+\beta}}\simeq1-\beta.
\]

故

\[
\Delta\nu_{\rm D}\simeq\nu_0[(1+\beta)-(1-\beta)]
=2\beta\nu_0.
\]

也就是

\[
\Delta\nu_{\rm D}\simeq2\nu_0\frac vc.
\]

代入 \(\nu_0=c/\lambda_0\) 还可直接化成

\[
\Delta\nu_{\rm D}\simeq\frac{2v}{\lambda_0}.
\]

数值为

\[
\begin{aligned}
\Delta\nu_{\rm D}
&\simeq\frac{2\times700}{486.1\times10^{-9}}\,\mathrm{Hz}\\
&\simeq2.88\times10^9\,\mathrm{Hz}.
\end{aligned}
\]

所以

\[
\boxed{
\Delta\nu_{\rm D}\simeq2.88\,\mathrm{GHz}
=2.88\times10^3\,\mathrm{MHz}
}.
\]

该值是从最快接近原子到最快远离原子的``总宽度''。单侧最大频移约为
\(1.44\,\mathrm{GHz}\)。

\textbf{（5）麦克斯韦分布下的方均根频偏}

图 7-2
右侧表示纵向速度的高斯分布在线性映射后仍为频偏高斯分布。低速纵向多普勒频移为

\[
\Delta\nu\simeq-\frac{v_x}{c}\nu_0.
\]

负号只表示朝向约定，对方均根没有影响。两边平方并取统计平均：

\[
\langle(\Delta\nu)^2\rangle
=\frac{\nu_0^2}{c^2}\langle v_x^2\rangle.
\]

一维麦克斯韦分布满足能量均分：

\[
\frac12m_H\langle v_x^2\rangle=\frac12k_BT.
\]

因此

\[
\langle v_x^2\rangle=\frac{k_BT}{m_H}.
\]

方均根频偏定义为

\[
\Delta\nu_{\rm rms}
=\sqrt{\langle(\Delta\nu)^2\rangle}.
\]

所以

\[
\begin{aligned}
\Delta\nu_{\rm rms}
&=\frac{\nu_0}{c}\sqrt{\langle v_x^2\rangle}\\
&=\frac{\nu_0}{c}\sqrt{\frac{k_BT}{m_H}}.
\end{aligned}
\]

最终

\[
\boxed{
\Delta\nu_{\rm rms}
=\frac{\nu_0}{c}\sqrt{\frac{k_BT}{m_H}}
}.
\]

由于 \(\nu_0/c=1/\lambda_0\)，也可写成

\[
\boxed{
\Delta\nu_{\rm rms}
=\frac1{\lambda_0}\sqrt{\frac{k_BT}{m_H}}
}.
\]

\textbf{两种线宽的数量级比较}

第三问自然线宽约 \(5.3\,\mathrm{MHz}\)，第四问限速模型给出的多普勒总宽约
\(2.88\,\mathrm{GHz}\)，后二者相差约

\[\frac{2.88\times10^9}{5.31\times10^6}\simeq5.4\times10^2.\]

这说明在给定速度尺度下，热运动多普勒展宽远大于有限寿命自然展宽。真实谱线还可能包含碰撞展宽、仪器展宽等，观测线型是多种机制的卷积。

\textbf{激发态编号核验}

能级数值为

\[E_2=-3.40\,\mathrm{eV},\qquad E_4=-0.850\,\mathrm{eV}.\]

差值

\[E_4-E_2=2.550\,\mathrm{eV},\]

所以绝不能把题目中的第 \(3\) 激发态误读成主量子数 \(n=3\)。

\newpage

\section{第 8 题}\label{ux7b2c-8-ux9898}

在光滑水平地面上有一个倒扣的箱子，内部长度为 \(l\)，质量为
\(M\)。箱内有一个质量为 \(m\)
的小球，初始时小球位于箱内右壁处，箱与小球均静止。小球与箱壁之间的碰撞均为完全弹性碰撞。自
\(t=0\) 起，对箱子施加恒定水平外力 \(F\)。

\begin{enumerate}
\def\labelenumi{\arabic{enumi}.}
\item
  求在

  \[t=5\sqrt{\frac{Ml}{F}}\]

  时，小球与箱壁已经发生的碰撞次数。
\item
  求该时刻小球相对地面的位移。
\end{enumerate}

\begin{figure}[H]
\centering
\includegraphics[width=0.826\linewidth,height=\textheight,keepaspectratio,alt={图 8-1　恒力推动的箱体与内部弹性碰撞小球。}]{../figures/question-08.pdf}
\caption*{图 8-1　恒力推动的箱体与内部弹性碰撞小球。}
\end{figure}

\subsection{一、题型定位}\label{ux4e00ux9898ux578bux5b9aux4f4d-7}

箱体在恒力下加速，小球相对箱壁作分段匀加速运动。换到箱体参考系后，两次碰撞间的加速度固定，碰撞条件可用无量纲时间统一表示。

\subsection{二、知识点}\label{ux4e8cux77e5ux8bc6ux70b9-7}

\begin{itemize}
\tightlist
\item
  相对坐标 \(\xi=x_m-x_M\)。
\item
  两碰撞间相对加速度 \(\ddot\xi=-F/M\)。
\item
  完全弹性碰撞使相对速度反向。
\item
  匀加速往返时间。
\item
  零相对速度接触与有冲量碰撞的区别。
\item
  系统质心加速度 \(F/(M+m)\)。
\item
  质心坐标与相对坐标的反变换。
\item
  分段时间与精确根式化简。
\end{itemize}

\subsection{三、物理图像}\label{ux4e09ux7269ux7406ux56feux50cf-7}

初始时小球在右壁。箱向右加速后，右壁向右离开小球，小球在地面系中暂时静止；箱的左壁逐渐追上小球。第一次真正碰撞发生在左壁。

左壁碰撞后，相对速度反向，小球相对箱向右运动。但相对加速度始终向左，所以它到达右壁时相对速度恰好降为零，随后立即向左返回。这个右端事件没有有限冲量，不计作碰撞。真正的碰撞只发生在左壁，时刻依次为
\(t_0,3t_0,5t_0,\ldots\)。

\subsection{四、解题入口}\label{ux56dbux89e3ux9898ux5165ux53e3-7}

令 \(x_M\) 为箱左壁位置，\(x_m\) 为小球位置，定义

\[\xi=x_m-x_M,\]

所以 \(0\le\xi\le l\)，且初始

\[\xi(0)=l,\qquad\dot\xi(0)=0.\]

两次碰撞之间

\[\ddot x_m=0,\qquad\ddot x_M=\frac FM,\]

故

\[\ddot\xi=-\frac FM.\]

先求从右端静止运动到左壁的时间 \(t_0\)。然后把目标时间写成
\(t/t_0\)，判断已经越过哪些奇数倍碰撞时刻。

\subsection{五、解析}\label{ux4e94ux89e3ux6790-7}

图 8-2 左侧以 \(\xi/l\) 对 \(t/t_0\) 画出小球在箱内的往返：曲线到达
\(\xi=0\) 的红点才是左壁碰撞，到达 \(\xi=l\)
时相对速度为零，只是转向。右侧依次列出相对运动、质心运动和坐标还原公式。第（1）---（3）步使用左图，第（4）---（6）步使用右侧三个框。

\begin{figure}[H]
\centering
\includegraphics[width=0.782\linewidth,height=\textheight,keepaspectratio,alt={图 8-2　相对坐标的抛物线往返、碰撞时刻与质心分解。}]{../figures/solution-08.pdf}
\caption*{图 8-2　相对坐标的抛物线往返、碰撞时刻与质心分解。}
\end{figure}

\textbf{（1）第一次碰撞时间} 按图 8-2 左侧的归一化坐标，初态位于
\(\xi/l=1\)，第一次红点位于 \(\xi=0\)。记

\[a=\frac FM.\]

相对运动初始条件为

\[\xi(0)=l,\qquad\dot\xi(0)=0,\]

相对加速度为

\[\ddot\xi=-a.\]

所以第一次碰撞前

\[
\xi(t)=l-\frac12at^2.
\]

第一次碰到左壁时 \(\xi=0\)。令该时刻为 \(t_0\)：

\[
0=l-\frac12at_0^2.
\]

于是

\[
\frac12at_0^2=l,
\]

\[
t_0^2=\frac{2l}{a}=\frac{2Ml}{F}.
\]

因此

\[
\boxed{t_0=\sqrt{\frac{2Ml}{F}}}.
\]

碰撞前相对速度

\[
\dot\xi(t_0^-)= -at_0=-\sqrt{2al}.
\]

\textbf{（2）完全弹性碰撞在相对运动中的作用}

图 8-2
左侧每个红点处曲线斜率反号，对应完全弹性碰撞使相对速度反向。一维完全弹性碰撞满足恢复系数

\[e=-\frac{\dot\xi^+}{\dot\xi^-}=1.\]

所以

\[
\dot\xi(t_0^+)=-\dot\xi(t_0^-)=\sqrt{2al}.
\]

碰撞后小球相对箱向右运动，但仍有恒定相对加速度
\(-a\)。从左壁出发，经过时间 \(\tau\)：

\[
\xi=\sqrt{2al}\,\tau-\frac12a\tau^2.
\]

到达右端 \(\xi=l\) 时：

\[
l=\sqrt{2al}\,\tau-\frac12a\tau^2.
\]

由于 \(\sqrt{2al}=at_0\) 且 \(l=at_0^2/2\)，代入后

\[
\frac12at_0^2=at_0\tau-\frac12a\tau^2.
\]

约去 \(a/2\)：

\[
t_0^2=2t_0\tau-\tau^2.
\]

即

\[
(\tau-t_0)^2=0.
\]

所以

\[\tau=t_0.\]

此时相对速度

\[
\dot\xi=\sqrt{2al}-a\tau=at_0-at_0=0.
\]

因此在绝对时刻
\(2t_0\)，小球到达右壁但相对速度为零，没有碰撞冲量；随后相对加速度使它向左返回。

从右端静止返回左端又需 \(t_0\)，故第二次真正碰撞发生在

\[t=3t_0.\]

此后过程周期性重复，左壁碰撞时刻为

\[
\boxed{t=(2n-1)t_0,\qquad n=1,2,3,\ldots}.
\]

\textbf{（3）目标时刻的碰撞次数}

图 8-2 已把左壁红点标在
\(t/t_0=1,3,5,\ldots\)，而偶数时刻只是右壁零相对速度转向。题给目标时刻

\[
t_*=5\sqrt{\frac{Ml}{F}}.
\]

与 \(t_0=\sqrt{2Ml/F}\) 比较：

\[
\frac{t_*}{t_0}
=\frac{5\sqrt{Ml/F}}{\sqrt{2Ml/F}}
=\frac5{\sqrt2}
\simeq3.536.
\]

在这一时刻之前的左壁碰撞时刻是

\[t_0,\qquad 3t_0.\]

下一次碰撞在 \(5t_0\)，尚未发生。\(2t_0\)
处的右壁事件相对速度为零，不算碰撞。因此

\[
\boxed{N=2}.
\]

\textbf{（4）目标时刻的质心位移}

转到图 8-2 右侧橙色框。把箱和小球看作整体，水平外力为 \(F\)，总质量为
\(M+m\)。质心加速度

\[
a_{\rm cm}=\frac{F}{M+m}.
\]

初始质心静止，所以目标时刻质心位移

\[
\Delta X_{\rm cm}=\frac12a_{\rm cm}t_*^2.
\]

而

\[
t_*^2=25\frac{Ml}{F}.
\]

故

\[
\begin{aligned}
\Delta X_{\rm cm}
&=\frac12\frac{F}{M+m}\cdot25\frac{Ml}{F}\\
&=\frac{25}{2}\frac{Ml}{M+m}.
\end{aligned}
\]

所以

\[
\boxed{
\Delta X_{\rm cm}=\frac{25Ml}{2(M+m)}
}.
\]

\textbf{（5）目标时刻小球在箱内的相对位置}

回到左侧曲线，第二次红点发生在 \(3t_0\)。目标时刻位于 \(3t_0\) 与
\(4t_0\) 之间，令从该次碰撞到目标时刻的时间为

\[
\tau=t_*-3t_0.
\]

写成 \(t_0\) 的倍数：

\[
\tau=\left(\frac5{\sqrt2}-3\right)t_0.
\]

定义无量纲量

\[
s=\frac{\tau}{t_0}=\frac5{\sqrt2}-3.
\]

第二次碰撞后，小球从左壁出发，相对初速度仍为

\[u_{\rm rel}=at_0.\]

所以

\[
\xi(t_*)=at_0\tau-\frac12a\tau^2.
\]

利用 \(at_0^2=2l\)：

\[
\begin{aligned}
\xi(t_*)
&=at_0^2s-\frac12at_0^2s^2\\
&=2ls-ls^2\\
&=l(2s-s^2).
\end{aligned}
\]

逐步化简 \(s\)：

\[
2s=\frac{10}{\sqrt2}-6=5\sqrt2-6,
\]

\[
\begin{aligned}
s^2
&=\left(\frac5{\sqrt2}-3\right)^2\\
&=\frac{25}{2}+9-\frac{30}{\sqrt2}\\
&=\frac{43}{2}-15\sqrt2.
\end{aligned}
\]

因此

\[
\begin{aligned}
2s-s^2
&=5\sqrt2-6-\left(\frac{43}{2}-15\sqrt2\right)\\
&=20\sqrt2-\frac{55}{2}.
\end{aligned}
\]

故

\[
\boxed{
\xi(t_*)=\left(20\sqrt2-\frac{55}{2}\right)l
}.
\]

数值上 \(20\sqrt2-27.5\simeq0.784\)，确实位于 \(0\) 与 \(1\) 之间。

\textbf{（6）由质心坐标还原小球对地位移}

最后使用图 8-2 右侧橙色框的反变换。以箱左壁位置为 \(x_M\)，有

\[x_m=x_M+\xi.\]

系统质心坐标

\[
X_{\rm cm}=\frac{Mx_M+mx_m}{M+m}.
\]

代入 \(x_m=x_M+\xi\)：

\[
\begin{aligned}
X_{\rm cm}
&=\frac{Mx_M+m(x_M+\xi)}{M+m}\\
&=x_M+\frac{m}{M+m}\xi.
\end{aligned}
\]

所以

\[x_M=X_{\rm cm}-\frac{m}{M+m}\xi.\]

于是

\[
\begin{aligned}
x_m
&=x_M+\xi\\
&=X_{\rm cm}-\frac{m}{M+m}\xi+\xi\\
&=X_{\rm cm}+\frac{M}{M+m}\xi.
\end{aligned}
\]

初始时 \(\xi_0=l\)，所以小球位移为

\[
\Delta x_m=\Delta X_{\rm cm}+\frac{M}{M+m}(\xi_* -l).
\]

代入前述结果：

\[
\Delta x_m
=\frac{25Ml}{2(M+m)}
+\frac{Ml}{M+m}\left(20\sqrt2-\frac{55}{2}-1\right).
\]

把括号中的常数合并：

\[
\begin{aligned}
\frac{25}{2}+20\sqrt2-\frac{55}{2}-1
&=20\sqrt2+\frac{25-55}{2}-1\\
&=20\sqrt2-15-1\\
&=20\sqrt2-16.
\end{aligned}
\]

因此

\[
\boxed{
\Delta x_m=(20\sqrt2-16)\frac{Ml}{M+m}
}.
\]

\textbf{为什么碰撞时相对速度只需反向}

设碰撞前箱与球速度分别为 \(V^-,v^-\)，碰撞后为
\(V^+,v^+\)。完全弹性碰撞的恢复系数条件为

\[v^+-V^+=-(v^--V^-).\]

左边和右边正是碰撞后的、碰撞前的相对速度，因此

\[\dot\xi^+=-\dot\xi^-.\]

无论 \(M,m\)
的具体值如何，这一相对速度反向关系都成立。质量只影响质心速度如何分配，而不影响相对坐标中的镜面反射规则。

\textbf{直接逐次碰撞法的核验}

第一次碰撞前，球静止，箱速度为

\[V^-_1=\frac FMt_0=\sqrt{\frac{2Fl}{M}}.\]

套用一维弹性碰撞公式后，可以验证

\[v_1^+-V_1^+=V_1^-,\]

即相对速度由 \(-V_1^-\) 变为
\(+V_1^-\)。但继续到第二次碰撞时，绝对速度表达会迅速变长；相对坐标法把所有质量比依赖都推迟到最后的质心还原，是本题更适合教学的主方法。

\end{document}
